% concatenated from AAFH_fractional_KBO.tex
 \documentclass[11pt,handout]{article}%
%Submission to Optimization Methods and Software. Same as the test case of Troeltzsch \cite{bib:TW04}}
\usepackage{geometry}                % See geometry.pdf to learn the layout options. There are lots.
\usepackage{tikz} 
\usetikzlibrary{positioning}
%\usepackage[draft]{graphicx}
\usepackage{graphicx}
\usepackage{amssymb}
%\usepackage{showkeys}
\usepackage{epstopdf}
\usepackage{subfigure}  % use for side-by-side figures
\usepackage[sans]{dsfont}
%\usepackage[applemac]{inputenc}
\usepackage[english]{babel}
\usepackage{latexsym}
\usepackage{mathrsfs}
%\usepackage{amscd}
\usepackage{graphicx}
\usepackage{color}
\usepackage{float}
\usepackage{mathtools}
%\usepackage{mathabx}
\frenchspacing
\usepackage{amsmath}
\usepackage{amsfonts}
\numberwithin{equation}{section}
\usepackage{enumerate}
\usepackage{amsthm}
%\usepackage[numbers]{natbib}
\usepackage[title]{appendix}

\usepackage[bookmarks=true,colorlinks=true,linkcolor={blue},urlcolor={blue}, citecolor={blue},pdfstartview={XYZ null null 1.22}]{hyperref}%
%\usepackage{slashbox}

%\usepackage{showkeys}

%%% For margin notes
%\newcommand{\note}[1]{%
%  \marginpar[{\color{red}{\raggedleft\small\sffamily #1\\}}]{%
	%  {\color{red}  {\raggedright\small\sffamily #1\\}}}}
	%
	%\usepackage{geometry}
	%\geometry{letterpaper}
	\providecommand{\keywords}[1]{\textbf{{Key words.}} #1}
	
	%    \graphicspath{{pics/}}
	
	%%%%%%%%%%%%%%%%%%%%%%%%%%%%%%%%%%%%%%%%%%%%%%%%%%%%%%%%%%%%%%%%%%%%%%%%%%%%%%%%%%%%%%%%%%%%%%%%%%%%%%%%%%%%%%%%%%%%%%%%%%%%%%%%%%%%%%%%%%
	\def\BigO{ {\mathcal {O}}}
	\def\bigM{ {\mathcal {M}}}
	\def\bigF{ {\mathcal {F}}}
	\def\B{\hat \beta}
	\def\f{\hat f}
	\def\g{\hat g}
	\def\Q{\hat Q}
	\def\real{\mathbb{R}}
	\newcommand{\R}{\mathbb R}
	\newcommand{\N}{\mathbb N}
	\newcommand{\dt}{\Delta t}
	%\renewcommand{\O}{\Omega}
	\newcommand{\Rd}{\mathbb{R}^d}
	\newcommand{\ball}{B}
	\newcommand{\abs}[1]{\left|#1\right|}
	\newcommand{\norm}[1]{\left\|#1\right\|}
	\def\be#1\ee{\begin{equation}#1\end{equation}}
	\newcommand{\fer}[1]{(\ref{#1})}
	\newtheorem{theorem}{Theorem}[section]
	%\newtheorem{corollary}{Corollary}
	%\newtheorem{lemma}[theorem]{Lemma}
	\newtheorem{proposition}{Proposition}
	%\newtheorem{conjecture}{Conjecture}
	\newtheorem{assumptions}{Assumptions}
	
	\theoremstyle{definition}
	\newtheorem{algorithm}{Algorithm}[section]
	%\newtheorem{definition}[theorem]{Definition}
	\newtheorem{remark}{Remark}
	\newcommand{\ep}{\varepsilon}
	\newcommand{\eps}[1]{{#1}_{\varepsilon}}
	\newcommand{\xx}{{\bf x}}
	\newcommand{\rr}{{\bf r}}
	\newcommand{\vvel}{{\bf v}}
	%\newcommand{\vv}{{\bf v}}
	\newcommand{\BB}{{\bf B}}
	\newcommand{\EE}{{\bf E}}
	\def\RR{\mathbb R}
	\def\e{\varepsilon}
	\def\cm{c_{\textrm{max}}}
	\def\ps{\p_{\infty}}
	\def\Ur{\textcolor{black}{V_r}}
	\def\Ua{\textcolor{black}{V_a}}
	\def\p{\rho}
	
	\def\N{\mathcal N}
	\def\T{\mathcal T}
	\def\L{\mathcal L}
	%%%%%%%%%%%%%%%%%%%%%%%%%%%%%%%%%%%%%%%%%%%%%%%%%%%%%%
	
	\title{Consensus-based optimization with $\alpha$-stable jump processes}
	
	
	\author{ 
Pedro  Aceves-Sanchez\thanks{Department of Mathematics, University of Arizona, Tucson, Arizona, USA 
	(pedroas@arizona.edu)},  \and Giacomo Albi\thanks{Department of Computer Science, University of Verona, ITALY (giacomo.albi@univr.it)}, \and F. Ferrarese \thanks{Department of Mathematics and Computer Science \& Center for Modeling, Computing and Statistics (CMCS), University of Ferrara, via Machiavelli 30, 44121 Ferrara, ITALY (federica.ferrarese@unife.it)}, \and Michael Herty \thanks{Institute for Geometry and Applied Mathematics (IGPM), RWTH Aachen University, GERMANY(herty@igpm.rwth-aachen.de) and
	Extraordinary Professor, Department of Mathematics and Applied Mathematics,  University of Pretoria, South Africa.}  
	}
	
	\date{}
	
	\begin{document}

\maketitle

% REQUIRED
\begin{abstract} 
	In this paper, we introduce a novel variant of the CBO method that incorporates jumps according to an $\alpha$-stable stochastic process in a kinetic framework. This extension gives rise to nonlocal stochastic effects, which improve the exploration capabilities of the method. We formulate the method at the particle level, detailing the corresponding stochastic dynamics and its asymptotic behavior. In particular, through a Fourier-based representation, we derive the associated fractional Fokker–Planck equation, which naturally accounts for the nonlocal diffusion behaviors induced by $\alpha$-stable processes. As a central result, we establish a rigorous convergence result for the proposed approach. Finally, we evaluate the performance of the method through a set of numerical experiments. The results demonstrate the effectiveness of the $\alpha$-stable jump process and emphasize its potential advantages over standard diffusion-based methods, particularly in complex optimization settings.
\end{abstract}

\section{Introduction}\label{sec:intro} 
Gradient-free optimization methods have attracted significant interest in recent years, especially in settings where the objective function is nonconvex, nonsmooth, noisy, or prohibitively expensive to evaluate \cite{kennedy1995particle, blum2003metaheuristics, borghi2024kinetic}. These approaches are particularly effective in high-dimensional and multi-modal landscapes, where gradient-based methods may fail due to the presence of multiple local minima, lack of regularity, or unreliable derivative information \cite{borghi2023adaptive,ferrarese2025localized}.
Classical examples of gradient-free optimization techniques include particle swarm optimization (PSO) \cite{kennedy1995particle, grassi2021particle}, genetic algorithms (GA) \cite{michalewicz1996genetic,mitchell1995genetic}, and more general evolutionary strategies, that relies on population-based heuristics inspired by biological or social behaviors \cite{blum2003metaheuristics}. In PSO, candidate solutions evolve by combining individual memory with social interactions, while GA employs selection, mutation, and recombination mechanisms to iteratively improve a population of individuals.
More recently, consensus-based optimization (CBO) has emerged as a mathematically grounded alternative within the class of population-based, gradient-free methods \cite{carrillo2018analytical, fornasier2020consensus, fornasier2021consensus}. In CBO, a swarm of agents evolves according to stochastic dynamics that drive the particles toward a weighted consensus point biased by the objective function, while random perturbations ensure sufficient exploration of the search space. Closely related are kinetic-based optimization (KBO) methods, which can be interpreted as the kinetic (mesoscopic) description of a time-discretized version of the CBO dynamics, typically obtained via the Euler–Maruyama scheme \cite{benfenati2022binary}. More precisely, given an objective function $\mathcal{E}:\RR^d\to \RR$, which has a unique global minimum at $x^\star$, the goal is to solve
\begin{equation}
	\min_x \mathcal{E}(x).
\end{equation}
Hence, in the classical KBO approach, a population of $N$ agents is considered, where each particle at position $x_i$, for any $i=1, \ldots ,N$, evolves according to the update rule
\begin{equation}\label{eq:KBO}
	x_i' = x_i + \nu   (\bar{x}_{\mathcal{E}}(t) - x_i) + \sigma  D(\bar{x}_{\mathcal{E}}(t), x_i)   z_i,
\end{equation}
with $x'_i$ denoting the post-interaction position. Here, $\nu$ and $\sigma$ are positive parameters controlling the relative strength of drift and diffusion, $z_i \sim \mathcal{N}(0,1)$, for any $i=1,\ldots, N$, represents a stochastic perturbation, and $D(\bar{x}_{\mathcal{E}}(t), x_i)$ is a diffusion matrix. The quantity $\bar{x}_{\mathcal{E}}(t)$ represents an estimate of the global minimizer, computed as a weighted average of the particle positions \cite{PTTM2017}. At the mesoscopic level, the evolution of the particle distribution $f(x,t)$ can be described according to
\begin{equation}\label{eq:boltzmann_kbo}
	\partial_t f(x,t) = \eta \, \mathcal{Q}(f)(x,t),
\end{equation}
which corresponds to a Boltzmann-type equation, where $\eta>0$ is the interaction frequency and the operator $\mathcal{Q}(\cdot
)$ accounts for the gain and loss of mass at position $x$ and time $t>0$, resulting from the interaction dynamics \eqref{eq:KBO}.
This kinetic formulation provides a mesoscopic description of the particle dynamics via a Boltzmann equation, which can be efficiently simulated using Direct Simulation Monte Carlo methods \cite{pareschi2013interacting, nanbu1980direct, pareschi2001introduction}. Furthermore, its long-time behavior characterizes the convergence of the particle system and, under suitable asymptotic regimes, yields Fokker--Planck approximations that enable a tractable convergence analysis. Extensions incorporating local interaction mechanisms between particles, as well as hybrid strategies combining KBO with genetic algorithms, have been proposed in \cite{albi2023kinetic, benfenati2022binary}. Finally, the hydrodynamic limits of fractional Fokker-Planck equations have been recently studied. See, for instance \cite{pasCesbron2019}.

The aim of this paper is to introduce a novel KBO model incorporating, in addition to classical Brownian diffusion, stochastic perturbations driven by an $\alpha$-stable L\'evy process \cite{applebaum2009levy,ken1999levy,kyprianou2022stable,meerschaert2019stochastic}. This combined framework enhances exploration by allowing agents to perform both local random walks and nonlocal jumps, effectively coupling diffusive and jump-driven dynamics. The proposed approach extends the models introduced in \cite{kalise2023consensus}, where jump-diffusion CBO dynamics were constructed by combining Wiener processes with compound Poisson noise. We also refer to \cite{borghi2025swarm} for an alternative approach to account for large jumps in CBO-dynamics based on BGK-type kinetic equations.

In particular, the standard jump-diffusion KBO model at the particle level reads
\begin{equation}\label{eq:jump_diffusion}
	x_i' = x_i + \nu(t)\big(\bar{x}_\mathcal{E}(t) - x_i\big) 
	+ \sigma(t) D\big(\bar{x}_\mathcal{E}(t), x_i\big) z_i 
	+ \gamma(t) D\big(\bar{x}_\mathcal{E}(t), x_i\big) \tilde{z}_i,
\end{equation}
for $i=1,\ldots,N$, where
\begin{equation}\label{eq:finite_poisson_measure}
	z_i \sim \mathcal{N}(0,1), 
	\qquad 
	\tilde{z}_i \sim \sum_{j=1}^{N^i(t)} w_j.
\end{equation}
Here, $N^i(t)$ are independent Poisson processes with intensity $\lambda$, and $w_j$ are i.i.d.\ $d$-dimensional random variables with zero mean and probability density $\rho(w)$, representing the jump amplitudes. The coefficients $\nu(t)$ and $\sigma(t)$ are continuous, positive, and non-decreasing functions such that $\nu(t)\to \nu>0$ and $\sigma(t)\to \sigma>0$ as $t\to\infty$, while $\gamma(t)$ is continuous, non-negative, and non-increasing with $\gamma(t)\to \gamma\geq 0$. The matrix $D\big(\bar{x}_\mathcal{E}(t), x\big)$ denotes the diffusion tensor, and $\bar{x}_\mathcal{E}(t)$ is the consensus point defined as a weighted average of the particle positions.
The jump component can be equivalently expressed in terms of a Poisson random measure as
\begin{equation}\label{eq:jump_diffusion_continuos}
	\tilde{z}_i \sim \int_{\mathbb{R}} w \, \mathcal{N}^i(dt,dw),
\end{equation}
where $\mathcal{N}^i(dt,dw)$ are independent Poisson random measures with intensity $\nu(dw)\,dt$, and $\nu(dw)$ is a finite L\'evy measure \cite{ken1999levy,applebaum2009levy}.
Several variants of this model can be considered. For instance, the Poisson intensity $\lambda$ may depend on time, with $\lambda(t)\to 0$ as $t\to\infty$, or the jump noise may be shared across particles, leading to simultaneous jumps with independent amplitudes. 

Here, the main objective is to formally derive the corresponding Fokker-Planck equation for the particle system, which naturally incorporates a fractional diffusion operator to account for the $\alpha$-stable jumps. A central challenge of these models is that $\alpha$-stable processes generally have infinite variance, which complicates both analytical treatment and interpretation of the concentration and convergence properties. Hence, a complete convergence proof will be provided, relying on the use of the Wasserstein $p$-distance between the evolving particle distribution and a Dirac measure centered at the global minimizer, inspired by the approach in \cite{fornasier2024consensus}.

The rest of the paper is organized as follows. In Section \ref{sec:particles_representation} we introduce the KBO method with $\alpha$-stable process, focusing on the particle evolution. In Section \ref{sec:FPE} we formally derive the corresponding fractional Fokker-Planck equation. In Section \ref{sec:convergence} we present a detailed convergence proof for the KBO method with $\alpha$ stable process. In Section \ref{sec:numerical_methods}, we introduce the numerical method used to simulate the particle dynamics, and in Section \ref{sec:numerical_experiments}, we propose different numerical experiments to validate the effectiveness of the method, and compare it to the classical KBO method. Finally, in Section \ref{sec:conclusions} we present some future research lines. 
\section{Kinetic framework for CBO with \texorpdfstring{$\alpha$}{alpha}-stable processes}\label{sec:particles_representation}
We introduce a KBO model in which stochastic exploration is driven by a symmetric $\alpha$-stable L\'evy process, extending the jump-diffusion formulation in \eqref{eq:jump_diffusion}. This choice allows particles to perform nonlocal jumps and enhances exploration beyond standard Brownian diffusion. In contrast to compound Poisson noise, $\alpha$-stable processes exhibit infinite activity and heavy-tailed increments, capturing long-range effects in the particle dynamics.
To this aim, we consider $f(x,t)
\geq 0$ and $x\in \mathbb{R}^d$ the distribution of particles in position $x$ at time $t
\geq 0$, assuming $f(x,
\cdot)$ a probability density function.
We suppose that each particle changes its position according to the following updating rules 
\begin{subequations}\label{eq:KBO_alfa_stable}
	\begin{align}
		x'   &= x + \nu (\bar{x}_\mathcal{E}(t) - x), \\
		x''  &= x' + \sigma D(\bar{x}_\mathcal{E}(t),x')\, z  + \gamma D(\bar{x}_\mathcal{E}(t),x')\, \tilde{z},
	\end{align}
\end{subequations}
where $\nu,\sigma,\gamma \geq 0$, $z \sim \mathcal{N}(0,1)$, and $\tilde{z} \sim \mathcal{L}_t^\alpha$, with $\mathcal{L}_t^\alpha$ denoting a symmetric $\alpha$-stable L\'evy process with $\alpha \in [1,2)$. Such process introduces an infinite-activity jump mechanism, in contrast to compound Poisson noise with discrete jumps. In particular, for $\alpha \in [1,2)$ the displacements follow a heavy-tailed distribution.

By the L\'evy representation theorem, the process $\mathcal{L}_t^\alpha$ admits the representation
\begin{equation}
	\mathcal{L}_t^\alpha = \int_{0}^t w \,\tilde{N}(ds,dw),
	\qquad 
	\tilde{N}(ds,dw) = N(ds,dw) - \nu(dw)\,ds,
\end{equation}
where $\tilde{N}(ds,dw)$ is a compensated Poisson random measure with L\'evy measure
\begin{equation*}
	\nu(dw) = C_{\alpha} |w|^{-d-\alpha}\,dw.
\end{equation*}
In \eqref{eq:KBO_alfa_stable}, $D(\bar{x}_\mathcal{E}(t),x)$ is a diffusion matrix, defined as either 
\[
D(\bar{x}_\mathcal{E}(t),x) = \vert \bar{x}_\mathcal{E}(t)-x \vert 
\textrm{Id}_d, 
\]
in the case of isotropic diffusion, where $\textrm{Id}_d$ is the identity matrix of size $d \times d$, or 
\[
D(\bar{x}_\mathcal{E}(t),x)  = diag\{(\bar{x}_\mathcal{E}(t)-x)_1,\ldots,(\bar{x}_\mathcal{E} (t)-x)_d\},
\]
in the case of anisotropic diffusion. Furthermore, the term $\bar{x}_\mathcal{E} (t)$ is the global estimate of the best position of the minimizer, which is computed as a convex combination of particle locations weighted by the cost function:
\begin{equation}\label{eq:x_tot}
	\bar{x}_\mathcal{E} (t) = \frac{\int_{\mathbb{R}^d}x e^{-\beta \mathcal{E}(x)}f(x,t)dx}{\int_{\mathbb{R}^d} e^{-\beta\mathcal{E}(x)}f(x,t)dx} \, .
\end{equation} 
The choice of the consensus point is made according to the Laplace principle \cite{dembo2009large, miller2006applied}, which yields
\begin{equation}\label{eq:laplace}
	\lim_{\beta \to \infty} 
	\left(
	-\frac{1}{\beta} \log \int_{\mathbb{R}^d} e^{-\beta \mathcal{E}(x)} f(x,t)\,dx
	\right)
	=
	\inf_{x \in \mathrm{supp}\, f(x,t)} \mathcal{E}(x).
\end{equation}

\paragraph{Boltzmann description.}
To describe the evolution of the particle density we rely on a kinetic approach, where $f(x,t)$ is governed by the Boltzmann-type equation
\begin{align}\label{eq:Boltz}
	\partial_t f(x,t) = \eta_1 \mathcal{Q}_1(f)(x,t) 
	+ \eta_2 \mathcal{Q}_2(f)(x,t),
\end{align}
with $\eta_i>0$, $i=1,2$, denote the interaction frequencies associated with \eqref{eq:KBO_alfa_stable}, namely the drift and diffusion mechanisms, respectively.
In weak form, for a given test function $\psi \in C_c^\infty(\mathbb{R}^d)$, the Boltzmann equation 
\eqref{eq:Boltz} reads
\begin{align}\label{eq:weak_Boltz_splitting}
	&\int_{\mathbb{R}^d} \partial_t f(x,t)\,\psi(x)\,dx
	= \eta_1 \int_{\mathbb{R}^d} \big[\psi(x') - \psi(x)\big] f(x,t)\,dx \cr
	&\quad + \eta_2 \int_{\mathbb{R}^{3d}} \big[\psi(x'') - \psi(x')\big] f(x',t)\, g_2(z)\,g_\alpha(\tilde z)\,dx\,dz\,d\tilde z.
	%&\qquad + \eta_3 \int_{\mathbb{R}^{2d}} \big[\psi(x''') - \psi(x)\big] f(x,t)\, g_\alpha(\tilde z)\,dx\,d\tilde z.
\end{align}
Here, the probability density $g_\alpha(\cdot
)$ encodes the distribution of the jump increments, and it is defined as the inverse Fourier transform of the characteristic function of a symmetric $\alpha$-stable distribution, namely,
\begin{equation*}
	g_\alpha(z) = \mathcal{F}^{-1}_{\kappa \to z}\big(\phi_\alpha(\kappa)\big)(z),
	\qquad 
	\phi_\alpha(\kappa) = \exp\big(-|\kappa|^\alpha\big),
\end{equation*}
where $\alpha \in (0,2]$. In particular, $\alpha=2$ corresponds to Gaussian noise, while $\alpha=1$ it corresponds to the Cauchy distribution. For the sake of completeness, we recall that the Fourier transform of a sufficiently regular function $f$ is given by
\[
\hat{f}(\xi) = \mathcal{F}(f)(\xi) := \int_{\mathbb{R}^d} e^{-\imath \xi \cdot x} f(x)\,dx,
\qquad \xi \in \mathbb{R}^d,
\]
with inverse transform
\[
\mathcal{F}^{-1}(h)(x) := \frac{1}{(2\pi)^d} \int_{\mathbb{R}^d} e^{\imath \xi \cdot x} h(\xi)\,d\xi.
\]
For notational simplicity, in what follows, we omit the dot product symbol in the exponential when no ambiguity arises.

Hence,  following \cite{pareschi2013interacting}, in \eqref{eq:weak_Boltz_splitting} we select the test function $\psi(x) = e^{-\imath \xi x}$, which corresponds to the Fourier representation of the Botlzmann-type equation 
\eqref{eq:Boltz} as
\begin{equation}\label{eq:weak_Boltz_fourier}
	\partial_t \hat{f}(\xi,t) 
	= \eta_1 \hat{\mathcal{Q}}_1(f)(\xi,t) 
	+ \eta_2 \hat{\mathcal{Q}}_2(f)(\xi,t),
	%+ \eta_3 \hat{\mathcal{Q}}_3(f)(\xi,t),
\end{equation}
where each term $\hat{\mathcal{Q}}_i(f)$, $i=1,2$ denotes the Fourier transform of the corresponding interaction operator.

In the following section we will focus on the formal derivation of a fractional Fokker-Planck model, describing the asymptotic property of the optimization process.

\subsection{Formal derivation of fractional Fokker-Planck}\label{sec:FPE} 
We now investigate the asymptotic behavior of \eqref{eq:weak_Boltz_fourier} under a suitable interaction scaling, analyzing separately the contribution of each operator in order to derive a Fokker--Planck-type approximation. Hence, we introduce the following scaling 
\begin{equation}\label{eq:grazing_scaling}
	\nu \to \nu \varepsilon, \qquad
	\sigma \to \sigma \sqrt{\varepsilon}, \qquad
	\gamma \to \gamma \varepsilon^{1/\alpha}, \qquad
	\eta_i \to {1}/{\varepsilon},
\end{equation}
which corresponds to a regime of frequent interactions with small increments, see for example
\cite{pareschi2013interacting} for further details. 
\paragraph{Scaled drift term.}
The first operator of \eqref{eq:weak_Boltz_splitting} write as follows
\begin{equation}\label{eq:operator_Q1_fourier}
	\hat{\mathcal{Q}}_1(f) 
	= \int_{\mathbb{R}^d} \big[e^{-\imath \xi x'} - e^{-\imath \xi x}\big] f(x,t)\,dx,
\end{equation}
with scaled interaction $x' = x + \varepsilon\nu(\bar{x}_\mathcal{E}(t)-x)$. Performing a Taylor expansion around $\varepsilon=0$, we obtain
\begin{equation}
	e^{-\imath \xi \nu( \bar{x}_\mathcal{E}(t) -x) \varepsilon}= 1 - \imath \xi \nu ( \bar{x}_\mathcal{E}(t) -x) \varepsilon + \mathcal{O}(\varepsilon^2),  %\frac{1}{2}\xi^2 \nu^2  \vert \bar{x}_\mathcal{E}(t) -x\vert^2 e^{-\imath \nu\xi (\bar{x}_\mathcal{E}(t)-x) \bar{\varepsilon}} \varepsilon^2,
\end{equation}
%with $\bar{\varepsilon} \in (0,\varepsilon)$. 
Considering the scaling \eqref{eq:grazing_scaling} for $\eta_1$  we have
\begin{equation}
	\eta_1 \hat{\mathcal{Q}}_1(f) = \frac{1}{\varepsilon} \left[ -\imath \xi \nu \varepsilon \int_{\RR^d} (\bar{x}_\mathcal{E}(t) -x)  e^{-\imath \xi x} f(x,t) dx  + \mathcal{O}(\varepsilon^2) \right].%R_1(\bar{\varepsilon}) \varepsilon^2 \right],   
\end{equation}
Retaining only the leading-order contribution, as $\varepsilon\to 0$ we formally have up to order $\mathcal{O}(\varepsilon)$
\begin{align}\label{eq:drift_fourier_final}
	\eta_1 \hat{\mathcal{Q}}_1(f)
	&= -\imath \xi \nu \int_{\mathbb{R}^d} (\bar{x}_\mathcal{E}(t) - x) e^{-\imath \xi x} f(x,t)\,dx\cr
	&\quad= -\nu\, \mathcal{F}\big(\nabla_x \cdot ((\bar{x}_\mathcal{E}(t)-x) f(x,t))\big).
\end{align}

\paragraph{Scaled diffusion term.}
Recalling \eqref{eq:weak_Boltz_splitting}, we write the second operator as
\begin{align}\label{eq:operator_Q2}
	\hat{\mathcal{Q}}_2(f)
	&= \int_{\mathbb{R}^{2d}} f(x',t) e^{-\imath \xi x'} 
	\int_{\mathbb{R}^d} \left(e^{-\imath \xi D(\bar{x}_\mathcal{E}(t),x') ( \sigma z + \gamma \tilde z)} - 1
	\right)g_2(z)g_\alpha(\tilde z) \,dz\,d \tilde z \, dx',\cr
	%\\&\quad - \int_{\mathbb{R}^d} f(x,t) e^{-\imath \xi x}\,dx \int_{\mathbb{R}^d} g_2(z)\,dz.
	&\quad= \int_{\mathbb{R}^d} e^{-\imath \xi x'} f(x',t)\, 
	\phi_2\big(\sigma D(\bar{x}_\mathcal{E}(t),x')\xi\big)\, \phi_\alpha\big(\gamma D(\bar{x}_\mathcal{E}(t),x')\xi\big)dx'
	- \hat{f}(\xi,t),
\end{align}
where we used that $\int_{\mathbb{R}^d} g_2(z)\,dz = 1$, $\int_{\mathbb{R}^d} g_\alpha(z)\,dz = 1$ and the Fourier representation of $g_2(\cdot)$ as
$\phi_2(\xi)=e^{-|\xi|^2}$ and of $g_\alpha(z)$ as $\phi_\alpha(\xi) = e^{-|\xi|^\alpha}$. Therefore, according to \eqref{eq:grazing_scaling} and using Taylor around $\varepsilon=0$, we have
\begin{equation}
	\begin{split}
		&\phi_2(\sigma \sqrt{\varepsilon} D(\bar{x}_\mathcal{E}(t),x') \xi)
		=1 -\sigma^2 D^2(\bar{x}_\mathcal{E}(t),x') \xi^2 \varepsilon + \mathcal{O}(\varepsilon^2),\\%\frac{1}{2} \sigma^4 \xi^4 D^4(\bar{x}_\mathcal{E}(t),x) e^{-\sigma^2 D^2(\bar{x}_\mathcal{E}(t),x) \xi^2 \bar{\varepsilon}_1} \varepsilon^2
		&\phi_\alpha(\gamma \varepsilon^{1/\alpha} \vert D(\bar{x}_\mathcal{E}(t),x')\vert^\alpha \vert \xi\vert^\alpha) 
		%&= e^{-\gamma^\alpha \vert D(\bar{x}_\mathcal{E}(t),x)\vert^\alpha \vert\xi\vert^\alpha \varepsilon} \cr
		= 1 -\gamma^\alpha \vert D (\bar{x}_\mathcal{E}(t),x') \vert^\alpha \vert \xi\vert^\alpha \varepsilon + \mathcal{O}(\varepsilon^2).
		%\cr
		%&\qquad\qquad+ \dfrac{1}{2}\gamma^{2\alpha} \vert \xi\vert^{2\alpha}\vert D  (\bar{x}_\mathcal{E}(t),x)\vert^{2\alpha} e^{-\gamma^\alpha \vert D (\bar{x}_\mathcal{E}(t),x)\vert^\alpha \vert\xi\vert^\alpha \bar{\varepsilon}_2} \varepsilon^2,
	\end{split}
\end{equation}
Therefore, scaling $\eta_2\rightarrow 1/\varepsilon$ we obtain,
\begin{align}
	\eta_2 \hat{\mathcal{Q}}_2(f) = &\frac{1}{\varepsilon} \Big[-\sigma^2 \xi^2 \varepsilon \int_{\RR^d} D^2(\bar{x}_\mathcal{E}(t),x') e^{-\imath \xi x'} f(x',t) dx'  \Big. \\
	& \Big. - \gamma^\alpha \vert \xi\vert^\alpha \varepsilon \int_{\RR^d} \vert D(\bar{x}_\mathcal{E}(t),x')\vert^\alpha e^{-\imath \xi x'} f(x',t) dx' + \mathcal{O}(\varepsilon^2) \Big]% \varepsilon^2 R_2(\bar{\varepsilon}_1) \Big],
\end{align}
Recalling that $x' = x + \mathcal{O}(\varepsilon)$, as $\varepsilon\to 0$, we formally have up to order $\mathcal{O}(\varepsilon)$
\begin{align}\label{eq:diffusion_fourier_final}
	\eta_2 \hat{\mathcal{Q}}_2(f)
	&= -\sigma^2 \xi^2 
	\int_{\mathbb{R}^d} D^2(\bar{x}_\mathcal{E}(t),x) e^{-\imath \xi x} f(x,t)\,dx\\ & \qquad - \gamma^\alpha \vert\xi\vert^\alpha 
	\int_{\mathbb{R}^d} \vert D(\bar{x}_\mathcal{E}(t),x)\vert^\alpha e^{-\imath \xi x} f(x,t)\,dx\cr 
	&=\sigma^2 \mathcal{F}\big(\Delta(D^2(\bar{x}_\mathcal{E}(t),x) f(x,t))\big) - \gamma^\alpha \vert \xi \vert^\alpha  \mathcal{F}\big(\vert D(\bar{x}_\mathcal{E}(t),x)\vert^\alpha f(x,t))\big).
\end{align}

Collecting the leading-order contributions in \eqref{eq:drift_fourier_final}, and \eqref{eq:diffusion_fourier_final}, we formally obtain the  following fractional Fokker--Planck equation
\begin{equation}\label{eq:fFP}
	\begin{aligned}
		\partial_t f(x,t) = 
		& - \nabla_x \cdot \big(\nu(\bar{x}_\mathcal{E}(t)-x) f(x,t)\big) + \sigma^2 \Delta\big(D^2(\bar{x}_\mathcal{E}(t),x) f(x,t)\big) \\
		& 
		\qquad - \gamma^{\alpha} (-\Delta)^{\alpha/2}\big(|D(\bar{x}_\mathcal{E}(t),x)|^\alpha f(x,t)\big),
	\end{aligned}
\end{equation}
where the fractional Laplacian encodes the nonlocal behavior induced by $\alpha$-stable jumps, and is defined as follows
\[
(-\Delta)^{\alpha/2} g(x) = \mathcal{F}^{-1}\big(|\xi|^\alpha \mathcal{F}(g)(\xi)\big).
\]

\section{Convergence to the global minimizer}\label{sec:convergence} 
The convergence analysis of consensus type global optimizers typically proceeds by first proving that a suitable measure of dispersion of the particle law, usually the variance,   decays in time exponentially, then showing that the particles mean converges to a limit point $\tilde{x}$, and finally arguing that this limit is in fact a global minimizer of the objective \cite{PTTM2017,benfenati2022binary}. 
For the kinetic KBO model driven by a symmetric $\alpha$-stable L\'evy process the same strategy remains natural but requires to be adapted. Indeed, for $\alpha<2$ the second moment is infinite, so one must replace variance estimates by bounds on $p$-th moment with $0<p<\alpha$ or by studying the behaviour of the Wasserstein $p$-distance between the distribution function and a Dirac centred in the minimization point, \cite{fornasier2024consensus}. Let $f(t,x) \in \mathcal{P}(\R\times \R^d)$ (space of probability measures in $\R^d$), $p \in [1, \infty)$ and $x^\star \in \R^d$, then the following identity holds
\begin{equation}\label{eq:Vp}
	W_p^p(f,\delta_{x^\star}) =   \int_{\R^d} |x - x^\star|^p f(t,x) dx := V_p(t),
\end{equation}
where with $\vert\cdot \vert$ we denote the Euclidean norm.
For any two probability measures $f_1,f_2 \in \mathcal{P} (\R^d)$, the Wasserstein distance $W_p$ between $f_1$ and $f_2$, for $p \in [1, \infty)$, is defined as

\[
W_p (f_1,f_2) = \inf_{\gamma \in \Gamma(f_1,f_2)} \bigg[ \int_{\R^d\times\R^d} |x - y|^p d \gamma(x,y) \bigg]^{1/p} \, ,
\]
where the set of transport plans from $f_1$ to $f_2$ is defined as

\[
\Gamma(f_1,f_2) := \big\{ \gamma \in \mathcal{P}(\R^d \times \R^d) : \pi^1_\# \gamma = f_1 \text{ and } \pi^2_\# \gamma = f_2 \big\}\, .
\]
Here $\pi^1_\# \gamma$ and $\pi^2_\# \gamma$ are the first and second marginals, respectively, of $\gamma$, \cite{ambrosio2008gradient}. The goal is to prove that $V_p(t)$ decreases exponentially fast, up to a prescribed accuracy $\varepsilon$, as $t \to \infty$, implying convergence in $W_p$ to the Dirac mass at the global minimizer. 

We now establish the main convergence result in the simplified setting with drift and pure $\alpha$-stable noise, i.e. with $\sigma = 0$ in \eqref{eq:fFP}. More specifically, we restrict our analysis to the case of isotropic diffusion. After introducing the required assumptions, we derive upper and lower bounds for $\frac{d}{dt}V_p(t)$ and prove that the mass around the minimizer remains positive for all times. These results allow us to show that the Wasserstein distance between $f$ and the Dirac mass at the minimizer decays exponentially fast as $t\to\infty$. The analysis relies on a quantitative version of the Laplace principle \cite{fornasier2024consensus}.
\begin{assumptions}\label{prop:assumptions} 
	Let us assume the following holds:
	\begin{itemize} 
		\item there exists $x^\star\in \R^d$ such that 
		\begin{equation*}
			\mathcal{E}(x^\star) = \inf_{x\in\R^d} \mathcal{E}(x) =: \underline{\mathcal{E}}, 
		\end{equation*}
		\item there exists $\mathcal{E}_\infty$, $R_0$, $\zeta>0$ and $\mu\in(0,\infty)$ such that
		\begin{equation}
			\begin{split}
				&\vert x-x^\star \vert \leq \frac{(\mathcal{E}(x)-\underline{\mathcal{E}})^\mu}{\zeta} \qquad \text{for all } x\in \mathcal{B}_{R_0}(x^\star),\\  
				&\mathcal{E}(x)-\underline{\mathcal{E}} > \mathcal{E}_\infty \qquad \text{for all } x\in \left( \mathcal{B}_{R_0}(x^\star)\right)^c. 
			\end{split}
		\end{equation} 
	\end{itemize}
\end{assumptions}
Then, the following holds.  
\begin{proposition}\label{prop:proposition1}
	For any $1<p<\alpha<2$, the fractional $p$-th moment in \eqref{eq:Vp} satisfies 
	\begin{equation}\label{eq:dtVp}
		\frac{d}{dt}V_p(t) \leq -C_{p,\alpha} V_p(t) + \nu p \vert \bar{x}_\mathcal{E}(t) - x^\star \vert (V_p(t))^{\frac{p-1}{p}} + \gamma^\alpha B_{p,\alpha} \vert \bar{x}_\mathcal{E}(t) - x^\star \vert^\alpha V_{p-\alpha}(t),
	\end{equation}
	where we recall $\vert \cdot \vert$ denotes the Euclidean norm. In \eqref{eq:dtVp} we define 
	\begin{equation*}
		V_{p-\alpha}(t) = \int_{\R^d\setminus\{x^\star\}} \vert x-x^\star \vert^{p-\alpha} f(x,t) dx,
	\end{equation*} and 
	\begin{equation}\label{eq:Cpalpha}
		C_{p,\alpha} = \nu p - \gamma^\alpha B_{p,\alpha}, \qquad \text{ with }	B_{p,\alpha} = \frac{2^\alpha \Gamma(\frac{d+p}{2}) \Gamma(\frac{\alpha-p}{2})}{\vert \Gamma (-\frac{p}{2})\vert + \Gamma(\frac{d+p-\alpha}{2})}>0,
	\end{equation}
	and with $\nu,\gamma>0$.
	Secondly, relying on Proposition 4.5 of \cite{fornasier2024consensus,fornasier2025pde}  (quantitative Laplace principle) and on Lemma 3.1 in \cite{harjulehto2018pointwise} the following estimates hold
	\begin{equation}\label{eq:estimates_Vp}
		\begin{split}
			&\vert \bar{x}_\mathcal{E}-x^\star\vert  \leq   \left[ C_r  + \left( \frac{e^{-\beta q}}{  \rho_t(\mathcal{B}_r(x^\star)) }\right)  (V_p(t))^{\frac{1}{p}}\right],\\
			&V_{p-\alpha}(t)\leq \alpha p \frac{\omega_d}{C(d)(p-\alpha)^2} (V_p(t))^{\frac{p-\alpha}{p}},  
		\end{split}
	\end{equation}
	where 
	\begin{equation}\label{eq:def_quantitative_laplace}
		\omega_d = \frac{\pi^{\frac{d}{2}}}{\Gamma(\frac{d}{2}+1)},\qquad C_r = \frac{(q+\mathcal{E}_r)^\mu}{\zeta}, \qquad \text{with }\quad  \mathcal{E}_r := \sup_{x\in\mathcal{B}_r(x^\star)} \mathcal{E}(x) \, , 
	\end{equation}
	$r > 0$, $q>0$ s.t. $q+\mathcal{E}_r\leq \mathcal{E}_\infty$, 
	$C(d)$ a constant dependent only on the dimension $d$, $\rho_t(\mathcal{B}_r(x^\star))$ denoting the mass in a ball of radius $r$ around the minimizer. 
\end{proposition}
\begin{proof}
	Consider the fractional Fokker-Planck equation in \eqref{eq:fFP} and set $\sigma=0$. Compute 
	\begin{equation}\label{eq:proof1} 
		\begin{split}
			\frac{d}{dt}V_p(t) = &\frac{d}{dt}\int_{\R^d} \vert x-x^\star \vert^p f(x,t) dx = \int_{\R^d} \vert x-x^\star \vert^p \partial_t f(x,t) dx = \\
			& = -\nu \int_{\R^d} \vert x-x^\star \vert^p \nabla_x\left(  (\bar{x}_\mathcal{E}(t)-x)f(t,x)\right)  dx +\\
			&\quad -\gamma^\alpha \int_{\R^d} \vert x-x^\star \vert^p (-\Delta)^{\alpha/2}\left(  \vert \bar{x}_\mathcal{E}(t)-x\vert^\alpha f(x,t)\right)  dx,
		\end{split} 
	\end{equation}
	where we have assumed $D(\bar{x}_\mathcal{E},x) = \vert \bar{x}_\mathcal{E}(t)-x\vert$. Let us now focus on the drift and diffusion terms separately. We start with the drift term and we integrate by parts to get 
	\begin{align*}
		&-\nu \int_{\R^d} \vert x-x^\star \vert^p \nabla_x\left(  (\bar{x}_\mathcal{E}(t)-x)f(x,t)\right)  dx 	=\\&= \nu \int_{\R^d} \nabla_x (\vert x-x^\star\vert^p) (\bar{x}_\mathcal{E}(t)-x) f(x,t) dx = \\ 
		& = \nu p \int_{\R^d} \vert x-x^\star \vert^{p-2} \langle x-x^\star, \bar{x}_\mathcal{E}(t)-x\rangle f(x,t) dx =\\ & = \nu p \int_{\R^d}  \vert x-x^\star \vert^{p-2} \langle x-x^\star, \bar{x}_\mathcal{E}(t)-x^\star + x^\star - x\rangle f(x,t) dx \leq \\  & \leq -\nu p \int_{\R^d} \vert x-x^\star \vert^{p-2} \vert x-x^\star \vert^2 f(x,t) dx +\\&+ \nu p \int_{\R^d} \vert x-x^\star \vert^{p-2} \langle \bar{x}_\mathcal{E}(t)-x^\star, x-x^\star\rangle f(x,t)dx \leq \\
		&\leq  -\nu p \int_{\R^d} \vert x-x^\star \vert^{p}   f(x,t) dx + \nu p ~ \vert \bar{x}_\mathcal{E}(t)-x^\star \vert \int_{\R^d} \vert x-x^\star \vert^{p-1} f(x,t) dx =\\ & = -\nu p V_p(t) + \nu p \vert \bar{x}_\mathcal{E}(t)-x^\star \vert V_{p-1}(t)   \leq -\nu p V_p(t) + \nu p  \vert \bar{x}_\mathcal{E}(t)-x^\star \vert \left( V_p(t)\right)^{\frac{p-1}{p}},
	\end{align*} 
	where $\langle \cdot,\cdot \rangle$ denotes the scalar product, and where we used Cauchy-Schwarz and H{\"o}lder inequalities. Now, by the definition of the consensus point, we get
	\begin{equation*}
		\vert \bar{x}_\mathcal{E}(t)-x^\star \vert = \Bigg \vert \frac{\int_{\R^d}x e^{-\beta \mathcal{E}(x)} f(x,t) dx}{\int_{\R^d}  e^{-\beta \mathcal{E}(x)} f(x,t) dx} -x^\star\Bigg \vert \leq   \int_{\R^d} \vert x-x^\star \vert ~ \frac{\omega_\beta(x)}{\vert \omega_\beta(x)\vert_{L^1(f)}}  f(x,t) dx,
	\end{equation*}
	where we denoted by $\omega_\beta(x) = e^{-\beta \mathcal{E}(x)}$.
	Following the steps in Proposition 4.5 of \cite{fornasier2024consensus}, we get by Markov's inequality 
	\begin{equation*}
		\vert \omega_\beta(x)\vert_{L^1(f)}\geq e^{-\beta \mathcal{E}_r} \rho_t(\mathcal{B}_r(x^\star)), 
	\end{equation*}
	being $\mathcal{E}_r$ defined as in \eqref{eq:def_quantitative_laplace} and $\rho_t(\mathcal{B}_r(x^\star))$ the mass in a ball of radius $r$ centred in the position of the global minimizer $x^\star$. Then, let $\tilde{r}>r>0$, we can write 
	\begin{equation*} 
		\begin{split}
			\vert \bar{x}_\mathcal{E}(t)-x^\star \vert & \leq    \int_{\mathcal{B}_{\tilde{r}}(x^\star)} \vert x-x^\star \vert ~  \frac{\omega_\beta(x) }{\vert \omega_\beta(x)\vert_{L^1(f)}}f(x,t) dx  +\\ &+   \int_{\left( \mathcal{B}_{\tilde{r}}(x^\star)\right) ^c} \vert x-x^\star \vert ~  \frac{\omega_\beta(x) }{\vert \omega_\beta(x)\vert_{L^1(f)}}f(x,t) dx:= T_r + T_r^c.
		\end{split}
	\end{equation*} 
	Now, on $\mathcal{B}_{\tilde{r}}(x^\star)$ we have $\vert x-x^\star\vert\leq \tilde{r}$ and so $T_r<\tilde{r}$. Additionally, 
	\begin{equation}
		\begin{split}
			T_r^c \leq &  \frac{e^{-\beta (\underline{\mathcal{E}}_r-\mathcal{E}_r)}}{  \rho_t(\mathcal{B}_r(x^\star))} \int_{\R^d} \vert x-x^\star \vert ~ f(x,t) dx \leq   \frac{e^{-\beta (\underline{\mathcal{E}}_r-\mathcal{E}_r)}}{  \rho_t(\mathcal{B}_r(x^\star))} \left( \int_{\R^d} \vert x-x^\star \vert^p ~ f(x,t) dx\right)^{\frac{1}{p}} ,
		\end{split}
	\end{equation}
	with \[\underline{\mathcal{E}}_r = \inf_{x \in \left( \mathcal{B}_{\tilde{r}}(x^\star)\right) ^c } \mathcal{E}(x),
	\] 
	and where we use Jensen inequality to have an estimate in terms of $V_p(t)$. 
	Let us choose $\tilde{r} = (q+\mathcal{E}_r)^\mu/\zeta$. Under Assumptions \ref{prop:assumptions} with $\underline{\mathcal{E}}=0$, $r\leq R_0$,  it can be easily verified that $\tilde{r}\geq r$, and $\tilde{r}\leq {\mathcal{E}_\infty^\mu}/{\zeta}. $
	Thus,
	\begin{equation*}
		\underline{\mathcal{E}}_r - \mathcal{E}_r \geq \min \{\mathcal{E}_\infty, (\zeta \tilde{r})^{\frac{1}{\mu}} \} - \mathcal{E}_r  = (\zeta \tilde{r})^{\frac{1}{\mu}} - \mathcal{E}_r = q,
	\end{equation*}
	since $(\zeta \tilde{r})^{\frac{1}{\mu}} = q+\mathcal{E}_r \leq \mathcal{E}_\infty$.  Finally,
	\begin{equation}\label{eq:estimate1} 
		\vert \bar{x}_\mathcal{E}(t) - x^\star \vert \leq  C_r + \frac{e^{-\beta q} }{\rho_t(\mathcal{B}_r(x^\star))}  (V_p(t))^{\frac{1}{p}},  
	\end{equation}
	with $C_r$ as in \eqref{eq:def_quantitative_laplace}. 
	Let us now focus on the fractional diffusion part in \eqref{eq:proof1}. As for the drift term, we begin by integrating by parts: 
	\begin{equation}\label{eq:proof_diffusion1} 
		\begin{split}
			&-\gamma^\alpha \int_{\R^d} \vert x-x^\star \vert^p (-\Delta)^{\alpha/2}\left(  \vert \bar{x}_\mathcal{E}(t)-x\vert^\alpha f(x,t)\right)  dx =\\ &= -\gamma^\alpha \int_{\R^d} (-\Delta)^{\alpha/2} ( \vert x-x^\star\vert^p) \vert \bar{x}_\mathcal{E}(t)-x\vert^\alpha f(x,t) dx. 
		\end{split}
	\end{equation}
	We then compute the fractional $\alpha$ derivative of the power $\vert x-x^\star\vert^p$ for $1<p<\alpha<2$ by relying on the  Fourier, as follows 
	\begin{equation}\label{eq:x_p}
		\begin{split}
			&(-\Delta)^{\alpha/2} ( \vert x-x^\star\vert^p) = \mathcal{F}^{-1} (\vert \xi \vert^\alpha \mathcal{F}(\vert x-x^\star \vert^p )) = \mathcal{F}^{-1}\left(\vert \xi \vert^\alpha \int_{\R^d} \vert y \vert^p e^{-\imath \xi (y+x^\star)} dx\right) \cr
			& = \mathcal{F}^{-1}\left( \vert \xi \vert^\alpha e^{-\imath \xi x^\star} \mathcal{F}(\vert y \vert^p) \right) =  \mathcal{F}^{-1}\left( D_p \vert \xi \vert^{-d-p+\alpha} e^{-\imath \xi x^\star}  \right),
		\end{split}
	\end{equation}
	where the constant
	\begin{equation*}
		D_p =  2^{p+d} \pi^{d/2} \frac{\Gamma(\frac{d+p}{2})}{\Gamma(-\frac{p}{2})},
	\end{equation*}
	comes from the  explicit computation of the Fourier transform of $\vert y \vert^p$ in dimension $d$. Now, we have
	\begin{equation}\label{eq:x_palpha}
		\mathcal{F}(\vert x-x^\star\vert^{p-\alpha}) = D_{p-\alpha} e^{-\imath \xi x^\star} \vert \xi \vert^{-d-p+\alpha}, 
	\end{equation}
	with 
	\begin{equation*}
		D_{p-\alpha} =  2^{p-\alpha+d} \pi^{d/2} \frac{\Gamma(\frac{d+p-\alpha}{2})}{\Gamma(-\frac{p-\alpha}{2})}.
	\end{equation*}
	Now by combining \eqref{eq:x_p} with \eqref{eq:x_palpha} we get
	\begin{equation*}
		\mathcal{F}((-\Delta)^{\alpha/2} (\vert x-x^\star \vert)) = \frac{D_p}{D_{p-\alpha}} \mathcal{F}(\vert x-x^\star \vert^{p-\alpha}),
	\end{equation*}
	and by taking the inverse Fourier transform, we get  
	\begin{equation}\label{eq:fractional_derivative_power}
		(-\Delta)^{\alpha/2} (\vert x-x^\star\vert^p) = - B_{p,\alpha} \vert x-x^\star \vert^{p-\alpha},
	\end{equation}
	with $B_{p,\alpha} $ as in \eqref{eq:Cpalpha}.
	Now by plugging \eqref{eq:fractional_derivative_power} into equation \eqref{eq:proof_diffusion1} we get 
	\begin{equation*}
		\begin{split}
			& \gamma^\alpha B_{p,\alpha} \int_{\R^d\setminus\{x^\star\}} \vert x-x^\star \vert^{p-\alpha} \vert \bar{x}_\mathcal{E}(t) -x \vert^\alpha f(x,t) dx =  \\
			& = \gamma^\alpha B_{p,\alpha}  \int_{\R^d\setminus\{x^\star\}} \vert x-x^\star \vert^{p-\alpha} \vert \bar{x}_\mathcal{E}(t) - x^\star + x^\star - x \vert^\alpha f(x,t) dx \leq \\
			& \leq \gamma^\alpha B_{p,\alpha}  \left[    \int_{\R^d\setminus\{x^\star\}}  \vert x-x^\star \vert^{p} f(x,t) dx +\right. \\ &\left. \qquad \qquad \quad +\vert \bar{x}_\mathcal{E}(t) - x^\star \vert^\alpha \int_{\R^d\setminus\{x^\star\}} \vert x-x^\star \vert^{p-\alpha} f(x,t) dx \right]=\\
			& = \gamma^\alpha B_{p,\alpha} \left( V_p(t) +\vert \bar{x}_\mathcal{E}(t) - x^\star \vert^\alpha ~ V_{p-\alpha}(t) \right).
		\end{split} 
	\end{equation*}
	Now,  the estimate of the term $\vert \bar{x}_\mathcal{E}(t)-x^\star \vert^\alpha$ follows straightforward by taking the power $\alpha$ in \eqref{eq:estimate1}:
	\begin{equation}
		\vert \bar{x}_\mathcal{E}(t) - x^\star \vert^\alpha  \leq \left(  C_r + \frac{e^{-\beta q} }{\rho_t(\mathcal{B}_r(x^\star))}  V_p^\frac{1}{p}(t)\right)^\alpha \leq 2^{\alpha-1}\left[  C_r^\alpha + \left( \frac{e^{-\beta   q} }{\rho_t(\mathcal{B}_r(x^\star))} \right)^\alpha (V_p(t))^\frac{\alpha}{p}\right].
	\end{equation}
	We conclude the proof by showing that fixing $r>0$, the second estimate in \eqref{eq:estimates_Vp} holds.
	We begin by decomposing $V_{p-\alpha}(t)$ into
	\begin{equation}
		V_{p-\alpha}(t)  = \int_{\mathcal{B}_r(x^\star)} \vert x-x^\star \vert^{p-\alpha} f(x,t) dx + \int_{(\mathcal{B}_r(x^\star))^c} \vert x-x^\star \vert^{p-\alpha} f(x,t) dx.
	\end{equation}
	Now, by Lemma 3.1 in \cite{harjulehto2018pointwise} we get for any $r>0$
	\begin{align}\label{eq:Vpalfa_1} 
		%\begin{split}
		&\int_{\mathcal{B}_r(x^\star)} \vert x-x^\star \vert^{p-\alpha} f(x,t) dx =  	\int_{\mathcal{B}_r(x^\star)} \vert x-x^\star \vert^{1-d} \vert x-x^\star \vert^{d-1+p-\alpha} f(x,t) dx 
		\cr
		& \qquad \qquad \leq \sup_{x \in \mathcal{B}_r(x^\star)} \vert x-x^\star \vert^{d-1+p-\alpha} \int_{\mathcal{B}_r(x^\star)} \vert x-x^\star\vert^{1-d}f(x,t) dx \leq \cr 
		&\qquad \qquad \leq  \frac{r^{d-1+p-\alpha} C(d) r}{\text{Vol}(\mathcal{B}_R(x^\star))} \int_{\mathcal{B}_R(x^\star)} f(x,t) dx \leq \frac{r^{d+p-\alpha} C(d) }{\omega_d r^d} = \frac{r^{p-\alpha} C(d) }{\omega_d},
		%\end{split}
	\end{align}
	where we recall that $\text{Vol}(\mathcal{B}_R(x^\star)) = \omega_d r^d$ with $\omega_d$ as in \eqref{eq:def_quantitative_laplace}.
	Secondly, 
	\begin{equation}\label{eq:Vpalfa_2}
		\int_{(\mathcal{B}_r(x^\star))^c} \vert x-x^\star \vert^{p-\alpha} f(x,t) dx\leq \frac{1}{r^{\alpha}} \int_{(\mathcal{B}_r(x^\star))^c}\vert x-x^\star\vert^p f(x,t) dx \leq  \frac{1}{r^{\alpha}} V_p(t).
	\end{equation}
	Now, optimizing over $r$ we get 
	\begin{equation*}
		r^\star = \left( \frac{\alpha \omega_d}{C(d)(p-\alpha)} V_p(t)\right)^{\frac{1}{p}},
	\end{equation*}
	and plugging it into \eqref{eq:Vpalfa_1}-\eqref{eq:Vpalfa_2} we get the second estimate in \eqref{eq:estimates_Vp}. 
\end{proof}

Before stating and proving the convergence theorem, we introduce the following auxiliary results. The first provides a lower bound for the quantity $\frac{d}{dt} V_p(t)$, while the second establishes a lower bound for the mass $\rho_t(\mathcal{B}(x^\star)$ for an arbitrarily small $r>0$.  
\begin{proposition}\label{prop:proposition2} 
	Under the assumptions of Proposition \ref{prop:proposition1}, the following estimate holds
	\begin{equation}\label{eq:lower_bound_Vp}
		\frac{d}{dt}V_p(t) \geq -(\nu p -\gamma^\alpha B_{p,\alpha}) V_p(t) - \nu p \vert \bar{x}_\mathcal{E}(t)-x^\star \vert (V_p(t))^{\frac{p-1}{p}},
	\end{equation}
	with $B_{p,\alpha}$ as in \eqref{eq:Cpalpha}, (see \cite{abatangelo2025gentle} for a detailed explanation on the computations of fractional Laplacians).
\end{proposition}
\begin{proof}
	The proof follows from Proposition \ref{prop:proposition1} by noting that by Cauchy-Schwarz and H{\"o}lder inequalities
	\begin{equation*}
		-\nu p \int_{\R^d} \vert x-x^\star \vert^{p-2} \langle \bar{x}_\mathcal{E}-x^\star, x-x^\star\rangle ~ f(x,t)dx \geq - \nu p \vert \bar{x}_\mathcal{E}(t)-x^\star \vert (V_p(t))^{\frac{p-1}{p}}.
	\end{equation*}
\end{proof}
\begin{proposition}\label{prop:mass_bounds}
	Let $\nu,\gamma>0$. Let $f$ be a weak solution of \eqref{eq:fFP}. For $r>0$, define the mollifier
	\begin{equation}\label{eq:mollifier}
		\phi_r(x) = \begin{cases}
			\exp \left( 1- \frac{r^2}{r^2 - \vert x-x^\star\vert^2} \right)  \qquad \text{if } \quad \vert x-x^\star \vert < r,\\
			0  \qquad  \qquad \qquad \qquad \qquad  \text{elsewhere}.
		\end{cases}
	\end{equation}
	Assume there exists $B>0$ such that $
	\vert \bar{x}_\mathcal{E}(t) - x^\star \vert \leq B$,
	for any $t \in [0,T]$, $T>0$ (follows from quantitative Laplace principle). Then, for any $t\in[0,T]$ we have 
	\begin{equation}\label{eq:mass_bound}
		\rho_t(\mathcal{B}_r(x^\star)) \geq  m_0 e^{-P t}:= m_\star , 
	\end{equation}
	with $m_0 = \rho_o(\mathcal{B}(x^\star))$ and 
	\begin{equation*}
		P = \nu \frac{C_{drift}(c)}{r} (B+cr) + \gamma^\alpha\frac{C_{frac}(c,\alpha)}{r^\alpha} (B+cr)^\alpha,
	\end{equation*}
	with $C_{drift}(c)>0$, $ C_{frac}(c,\alpha)>0$, for any $c\in (0,1)$.  
\end{proposition}
\begin{proof}
	By definition of the mollifier in \eqref{eq:mollifier} we have $0\leq \phi_r(x)\leq 1$ and $supp (\phi_r(x)) = \mathcal{B}_r(x^\star)$. Then 
	\begin{equation}\label{eq:mass_phi_relation}
		\rho_t(\mathcal{B}_r(x^\star)) \geq \int_{\RR^d} \phi_r(x) f(x,t) dx. 
	\end{equation}
	The idea is to find a lower bound for the right-hand side of \eqref{eq:mass_phi_relation}. Hence, we consider the weak form of \eqref{eq:fFP}
	\begin{equation*}
		\frac{d}{dt} \int_{\RR^d} \phi_r(x) f(x,t) dx = I_{drift}(t) + I_{frac}(t),
	\end{equation*}
	with
	\begin{equation*}
		\begin{split}
			I_{drift}(t) = & \nu \int_{\RR^d}  (\bar{x}_\mathcal{E}(t)-x) \nabla_x \phi_r(x) f(x,t) dx := \int_{\RR^d} T_1(x) f(x,t) dx \\
			I_{frac}(t) = & -\gamma^\alpha \int_{\RR^d} \vert \bar{x}_\mathcal{E}(t)-x\vert^\alpha (-\Delta)^{\alpha/2} \phi_r(x) f(x,t) dx := \int_{\RR^d} T_2(x) f(x,t) dx.
		\end{split}
	\end{equation*}
	Now we proceed by splitting the domain $\RR^d$ into different regions, depending on the sign of $T_1(x)$ and $T_2(x)$. 
	We begin by analyzing the drift term. We recall that 
	\begin{equation*}
		\nabla \phi_r(x) = \begin{cases}
			\frac{-2 r^2 (x-x^\star) }{(r^2 - \vert x-x^\star\vert^2)^2} \phi_r(x)  \qquad \text{if } \quad \vert x-x^\star \vert < r,\\
			0  \qquad  \qquad \qquad \qquad \qquad  \text{else}.
		\end{cases}
	\end{equation*}
	Hence, we get
	\begin{equation*}
		\text{sign}(T_1(x)) = - \text{sign}( \langle \bar{x}_\mathcal{E}(t)-x, x-x^\star\rangle). 
	\end{equation*}
	We define 
	\begin{equation*}
		\Omega_r = \{ x\in \RR^d : \vert x-x^\star\vert < r\},
	\end{equation*}
	and for any $c\in (0,1)$
	\begin{equation}
		\begin{split}
			& K_1 = \Omega_r^c \cup \{ x\in \Omega_r : \langle \bar{x}_\mathcal{E}(t)-x, x-x^\star\rangle \leq 0 \quad \text{and} \quad c r<\vert x-x^\star \vert < r \},\\
			& K_2 = \{ x\in \Omega_r : \langle \bar{x}_\mathcal{E}(t)-x, x-x^\star\rangle > 0 \quad \text{and} \quad \vert x-x^\star \vert \leq c r\}.
		\end{split}
	\end{equation}
	We note that $K_1 \cup K_2 = \RR^d$ and $K_1\cap K_2 = \emptyset$. 
	\paragraph{Case $K_1$.} 
	If $x\in \Omega_r$ and $\langle \bar{x}_\mathcal{E}(t)-x, x-x^\star\rangle\leq 0$ then 
	\begin{equation*}
		T_1(x) = \frac{-2 r^2 \nu}{(r^2 - \vert x-x^\star \vert^2)^2} \phi_r(x) \langle \bar{x}_\mathcal{E}(t)-x, x-x^\star\rangle \geq 0.
	\end{equation*}
	If $x \not \in \Omega_r$, then $\phi_r(x) = 0$ and $\nabla \phi_r(x)=0$, so $T_1(x)=0$. 
	Therefore, this yields
	\begin{equation*}
		\int_{K_1\cap \Omega_r} T_1(x) f(x,t) dx \geq 0. 
	\end{equation*}
	\paragraph{Case $K_2$.} If $x\in K_2$ then $\langle \bar{x}_\mathcal{E}(t)-x, x-x^\star\rangle> 0$ and $\vert x-x^\star \vert \leq cr$ for any $c\in (0,1)$.
	Now,
	\begin{equation*}
		\begin{split}
			T_1(x) = & -2 r^2 \nu \frac{\langle \bar{x}_\mathcal{E}(t)-x, x-x^\star\rangle}{(r^2-\vert x-x^\star\vert ^2)^2}\phi_r(x) \geq -2 r^2 \nu \frac{\vert \bar{x}_\mathcal{E}(t)-x\vert \vert x-x^\star\vert}{(r^2-\vert x-x^\star\vert ^2)^2} \phi_r(x) \geq \\ 
			& \geq -2 r^2 \nu \frac{\left( \vert \bar{x}_\mathcal{E}(t)-x^\star\vert + \vert x^\star - x \vert \right) \vert x-x^\star\vert}{(r^2-\vert x-x^\star\vert ^2)^2} \phi_r(x) \geq \frac{-2r^2\nu (B+cr) cr}{(r^2-c^2r^2)^2}\phi_r(x) = \\
			= & -\nu \frac{C_{drift}(c)}{r} (B+cr) \phi_r(x) : = -p_1 \phi_r(x),
		\end{split}
	\end{equation*}
	with $C_{drift}(c) = 2c/(1-c^2)^2$. Then, 
	\begin{equation}\label{eq:I_drift}
		I_{drift} = \int_{\RR^d} T_1(x) f(x,t) dx \geq \int_{K_2} T_1(x) f(x,t)dx \geq -p_1 \int_{\RR^d} \phi_r(x) f(x,t).
	\end{equation}
	We now focus on $I_{frac}$. 
	In analogy with the drift term, we split the domain into $\Omega_r^c$ and 
	\begin{align*}
		\Omega_{in} = \{x \in \Omega_r : \vert x-x^\star \vert \leq cr\}, \qquad
		\Omega_{bd} = \{x \in \Omega_r : cr< \vert x-x^\star \vert <r\}. 
	\end{align*}
	\paragraph{Case $\Omega_r^c$.} In $\Omega_r^c$, $\phi_r(x)=0$. Using the definition of fractional Laplacian 
	\begin{equation*}
		(-\Delta)^{\alpha/2} \phi_r(x) = c_\alpha PV \int_{\RR_d} \frac{\phi_r(x)-\phi_r(y)}{\vert x-y\vert^{1+\alpha}} dy = c_\alpha PV \int_{\RR_d} \frac{0-\phi_r(y)}{\vert x-y\vert^{1+\alpha}} dy \leq 0,
	\end{equation*}
	for $c_\alpha>0$
	\paragraph{Case $\Omega_{in}$.} 
	Define 
	\begin{equation}\label{eq:phi_1} 
		\phi_1(u) =   \begin{cases}
			\exp \left( 1- \frac{1}{1 - u^2} \right)  \qquad \text{if } \quad \vert u \vert < 1,\\
			0  \qquad  \qquad \qquad \qquad  ~ \text{else}. 
		\end{cases}  
	\end{equation}
	Note that $\phi_r(x) = \phi_1\left( \frac{x-x^\star}{r}\right)$. For any $x\in \RR^d$, set $u=\frac{x-x^\star}{r} $:
	\begin{equation}\label{eq:relation_phi1_phi_r} 
		\begin{split}
			(-\Delta)^{\alpha/2} \phi_r(x) =&   c_\alpha PV \int_{\RR_d} \frac{\phi_1(u)-\phi_1(\frac{y-x^\star}{r})}{\vert x-y\vert^{1+\alpha}} dy = \\ & = c_\alpha PV \int_{\RR_d} \frac{\phi_1(u)-\phi_1(z)}{(r\vert u-z\vert)^{1+\alpha}} r dz  = r^{-\alpha}  (-\Delta)^{\alpha/2} \phi_1(x).
		\end{split}
	\end{equation}
	If $x\in \Omega_{in}$ then $\vert x-x^\star\vert\leq cr$, which means $\vert ur\vert \leq cr$ that is $\vert u \vert \leq c$. Then on the compact set $[-c,c]^d \in (-1,1)^d$, $\phi_1\in \mathcal{C}_c^\infty$ and strictly positive. Hence it $\phi_1(u)$ attains a minimum $m(c)>0$, and $((-\Delta)^{\alpha/2}\phi_1(x))$  attains a maximum $M(\alpha,c)<+\infty$  respectively. 
	Then 
	\begin{equation}
		(-\Delta)^{\alpha/2} \phi_r(x) =  r^{-\alpha}  (-\Delta)^{\alpha/2} \phi_1(x) \leq  r^{-\alpha} \frac{M(\alpha,c)}{m(c)} m(c) \leq   r^{-\alpha} \frac{M(\alpha,c)}{m(c)}\phi_r(x), 
	\end{equation}
	where we recall that $\phi_1(u) = \phi_r(x)$ by definition.  
	\paragraph{Case $\Omega_{bd}$.} 
	Fix $d\in\mathbb{N}$ and $\alpha\in(1,2)$. Define the radial bump function $\phi_1(x)$ as in \eqref{eq:phi_1}. 	 
	Fix $\alpha\in(1,2)$ and consider $x$ such that $\eta<|x|<1$, for some $\eta\in(3/4,1)$ to be chosen later. Using the symmetric representation of the fractional Laplacian, we write
	\begin{equation*}
		(-\Delta)^{\alpha/2}\phi_1(x)
		= c_{d,\alpha}\int_{\mathbb{R}^d}\frac{2\phi_1(x)-\phi_1(x+h)-\phi_1(x-h)}{|h|^{d+\alpha}}\,dh,
	\end{equation*}
	so that it is enough to prove that the integral is negative.
	Define 
	\[
	m_0 := \min_{|y|\le 1/2}\phi_1(y)>0, 
	\qquad 
	\varepsilon_\eta := \sup_{\eta\le |x|<1}\phi_1(x)=\phi_1(\eta),
	\qquad 
	d_\eta := m_0-2\varepsilon_\eta.
	\]
	For $\eta$ sufficiently close to $1$, one has $d_\eta>0$.
	
	We split the integral into three regions. First, consider the translated set 
	\[ 
	S_x := x - \mathcal{B}_{1/2}(0) = \{ h = x-y : y\in \mathcal{B}_{1/2}(0)\}.
	\]   For $h\in S_x$ we have $\phi_1(x-h)\ge m_0$ and $\phi_1(x+h)=0$, hence
	\[
	2\phi_1(x)-\phi_1(x+h)-\phi_1(x-h) \le 2 \varepsilon_\eta-0-m_0 = -d_\eta<0.
	\]
	Moreover, $|h|\le 3/2$ on $S_x$, and since $\abs{S_x}=\omega_d 2^{-d+1}$, we conclude
	\begin{equation}\label{eq:contribution_1}
		\int_{S_x}\frac{2\phi_1(x)-\phi_1(x+h)-\phi_1(x-h)}{|h|^{d+\alpha}}\,dh
		\le - d_\eta\,\omega_d\,2^{-d+1}\Bigl(\frac{2}{3}\Bigr)^{d+\alpha}.
	\end{equation}
	Next, for $|h|<\rho$ with $\rho\le \eta-1/2$, a Taylor expansion yields
	\[
	|2\phi_1(x)-\phi_1(x+h)-\phi_1(x-h)| \le M_2 |h|^2,
	\]
	so that
	\begin{equation*}
		\left| \int_{|h|<\rho}\frac{2\phi_1(x)-\phi_1(x+h)-\phi_1(x-h)}{|h|^{d+\alpha}}\,dh\right| \leq M_2 \int_{|h|<\rho} \vert h \vert^{2-d-\alpha} dh  
		= M_2\,\omega_d\,\frac{\rho^{2-\alpha}}{2-\alpha}.
	\end{equation*}
	and in particular 
	\begin{equation}\label{eq:contribution_2}
		\int_{|h|<\rho}\frac{2\phi_1(x)-\phi_1(x+h)-\phi_1(x-h)}{|h|^{d+\alpha}}\,dh
		\le M_2\,\omega_d\,\frac{\rho^{2-\alpha}}{2-\alpha}.
	\end{equation}
	Finally, using $\phi_1\ge 0$, for $|h|\ge \rho$ we obtain
	\begin{equation}\label{eq:contribution_3}
		\int_{|h|\ge \rho}\frac{2\phi_1(x)-\phi_1(x+h)-\phi_1(x-h)}{|h|^{d+\alpha}}\,dh
		\le 2\varepsilon_\eta \int_{\abs{h}\geq \rho} \abs{h}^{-(d+\alpha)} dh = 2\varepsilon_\eta\,\omega_d\,\frac{\rho^{-\alpha}}{\alpha} .
	\end{equation}
	Combining the three contributions in \eqref{eq:contribution_1}-\eqref{eq:contribution_2}-\eqref{eq:contribution_3}, we get
	\begin{equation*}
		(-\Delta)^{\alpha/2}\phi_1(x)
		\le c_{d,\alpha}\,\omega_d\Bigg[
		- d_\eta\,2^{-d+1}\Bigl(\frac{2}{3}\Bigr)^{d+\alpha}
		+ \frac{M_2}{2-\alpha}\rho^{2-\alpha}
		+ \frac{2\varepsilon_\eta}{\alpha}\rho^{-\alpha}
		\Bigg].
	\end{equation*}
	Optimizing in $\rho>0$ gives the choice $\rho_*=\sqrt{2\varepsilon_\eta/M_2}$, for which
	\[
	\frac{M_2}{2-\alpha}\rho_*^{2-\alpha}
	+ \frac{2\varepsilon_\eta}{\alpha}\rho_*^{-\alpha}
	= \frac{2}{\alpha(2-\alpha)}(2\varepsilon_\eta)^{\frac{2-\alpha}{2}} M_2^{\frac{\alpha}{2}}.
	\]
	If $\rho_*\le \eta-1/2$, we deduce
	\begin{equation*}
		(-\Delta)^{\alpha/2}\phi_1(x)
		\le c_{d,\alpha}\,\omega_d\Bigg[
		- d_\eta\,2^{-d+1}\Bigl(\frac{2}{3}\Bigr)^{d+\alpha}
		+ \frac{2}{\alpha(2-\alpha)}(2\varepsilon_\eta)^{\frac{2-\alpha}{2}} M_2^{\frac{\alpha}{2}}
		\Bigg].
	\end{equation*}
	Thus, choosing $\eta$ sufficiently close to $1$ so that
	\begin{equation*}
		d_\eta\,2^{-d+1}\Bigl(\frac{2}{3}\Bigr)^{d+\alpha}
		>
		\frac{2}{\alpha(2-\alpha)}(2\varepsilon_\eta)^{\frac{2-\alpha}{2}} M_2^{\frac{\alpha}{2}},
	\end{equation*}
	we conclude that
	$(-\Delta)^{\alpha/2}\phi_1(x)<0$, for all  $\eta<|x|<1.
	$ 
	Now, since $(-\Delta)^{\alpha/2}\phi_r(x) = r^{-\alpha} (-\Delta)^{\alpha/2}\phi_1(x)$ as proved \eqref{eq:relation_phi1_phi_r}, we get that 
	\[
	(-\Delta)^{\alpha/2}\phi_r(x)<0 \qquad \text{for all } \eta<|x|<1.
	\]
	Therefore, gathering the previous estimates yields
	\begin{equation}\label{eq:I_frac}
		\begin{split}
			I_{frac} = &-\gamma^\alpha \int_{\RR^d} \vert \bar{x}_\mathcal{E} -x \vert^\alpha   (-\Delta)^{\alpha/2} \phi_r(x) f(x,t) dx \\
			& \geq -\gamma^\alpha \int_{\Omega_{in}} \vert \bar{x}_\mathcal{E} -x \vert^\alpha   (-\Delta)^{\alpha/2} \phi_r(x) f(x,t) dx \\
			& \geq -\gamma^\alpha (B+cr)^\alpha r^{-\alpha} C_{frac}(c,\alpha)  \int_{\RR^d} \phi_r(x) f(x,t) dx,% :=\\
			%&:= - p_2  \int_{\RR^d} \phi_r(x) f(x,t) dx,
		\end{split}
	\end{equation}
	with $C_{frac}(c,\alpha) =  M(\alpha,c)/m(c)$.
	Finally, by plugging together the drift and fractional diffusion parts \eqref{eq:I_drift}-\eqref{eq:I_frac} and by applying Gr{\"o}nwall inequality we get the desired mass bound in \eqref{eq:bound_mass}. 
\end{proof}
We can now proceed by stating and proving the convergence theorem, in analogy with Theorem 3.7 of \cite{fornasier2024consensus}. 
\begin{theorem}\label{prop:convergence}
	Let $\mathcal{E} \in C(\R^d)$ satisfying Assumption \ref{prop:assumptions}. Let $f_0$ be such that $x^\star \in supp(f_0)$. Fix $\varepsilon \in (0,V_p(0))$, $\theta \in (0,1)$, choose the parameters $\nu,\gamma$ such that $\nu p > \gamma^\alpha B_{p,\alpha}$ with $B_{p,\alpha}$ defined as in \eqref{eq:Cpalpha}. Define 
	\begin{equation}\label{eq:T_star} 
		T_\star = \frac{1}{C_{p,\alpha}(1-\theta)} \log\left( \frac{V_p(0)}{\varepsilon}\right), \qquad \text{with } C_{p,\alpha} \text{ as in \eqref{eq:Cpalpha}}.
	\end{equation}
	Assume there exist $m_0>0$ and $P>0$ s.t. for any $r>0$,
	\begin{equation}\label{eq:bound_mass}
		\rho_t(B_r(x^\star)) \geq m_0 e^{-Pt} :=m_\star,
	\end{equation}
	on $[0,T_\star]$, as given by Proposition \ref{prop:mass_bounds}. Then, there exists $\beta_0>0$ such that for any $\beta>\beta_0$ we have
	\begin{equation}
		V_p(T) = \varepsilon, \qquad \text{with } T\in \left[ \frac{1-\theta}{1+\frac{\theta}{2}} T_\star, T_\star \right].
	\end{equation}
	Furthermore, on the time interval $[0,T]$, $V_p(t)$ decays at least exponentially fast:
	\begin{equation}
		W_p^p(f,\delta_{x_\star}) = V_p(t) \leq V_p(0) \exp(-(1-\theta)C_{p,\alpha}t).
	\end{equation}
\end{theorem}

\begin{proof}
	Without loss of generality, we assume $\underline{\mathcal{E}}=0$. Let us choose the parameter $\beta$ as follows 
	\begin{equation}\label{eq:beta}
		\beta > \beta_0 := \frac{1}{q_\varepsilon} \log\left( \frac{2}{ m_\star c}\right), \qquad \text{with } 	q_\varepsilon = \frac{1}{2} \min \left\lbrace \left( \frac{\zeta c \varepsilon^{1/p}}{2} \right)^{\frac{1}{\mu}}, \mathcal{E}_\infty \right\rbrace 
	\end{equation}
	$m_\star$ as in \eqref{eq:bound_mass} 
	and 
	\begin{equation}\label{eq:c}
		c = \min \left\lbrace \frac{\theta}{2} \frac{\nu p -\gamma^\alpha B_{p,\alpha}}{\nu p},\left( \frac{\theta}{2} \frac{\nu p -\gamma^\alpha B_{p,\alpha}}{\gamma^\alpha B_{p,\alpha}G(d,p,\alpha) }\right)^{\frac{1}{\alpha}}\right\rbrace, 
	\end{equation}
	being 
	\begin{equation}\label{eq:G}
		G(d,p,\alpha) = \alpha p \frac{\omega_d}{C(d)(p-\alpha)^2},
	\end{equation}
	with $\omega_d$ and $C(d)$ defined as in Proposition \ref{prop:proposition1}. In addition, we introduce 
	\begin{equation}
		r_\varepsilon = \max_{r\in[0,R_0]} \left\lbrace \max_{x\in \mathcal{B}_r(x_\star)} \mathcal{E}(x) \leq q_\varepsilon \right\rbrace. 
	\end{equation}
	Let us now define the time horizon $T_\beta>0$ by 
	\begin{equation}
		T_\beta :=\sup \left\lbrace  t>0 : V_p(t')>\varepsilon \quad \text{and} \quad  \vert \bar{x}_\mathcal{E}(t)-x^\star \vert < K(t')~ \forall t'\in [0,t]\right\rbrace, 
	\end{equation}
	with $K(t) := c (V_p(t))^{\frac{1}{p}}$ and $c$ as in \eqref{eq:c}. \\
	Our aim is now to show that 
	\begin{equation}\label{eq:proof_goal}
		V_p(T_\beta) = \varepsilon, \quad \text{with} \quad T_\beta \in \left[ \frac{1-\theta}{1+\frac{\theta}{2}} T_\star, T_\star \right],
	\end{equation}
	and that we have at least exponential decay of $V_p(t)$ until time $T_\beta$, that is until accuracy $\varepsilon$ is reached. \\
	First we ensure $T_\beta>0$. This follows from the definition since $V_p(0)>\varepsilon$  and $\vert \bar{x}_\mathcal{E}(0)-x^\star \vert < K(0)$. Indeed,
	\begin{equation*}
		\begin{split}
			\vert \bar{x}_\mathcal{E}(0)-x^\star \vert \leq & 2 \left[ C_{r_\varepsilon} +   \frac{e^{-\beta q_\varepsilon}}{\rho_t(\mathcal{B}_{r_\varepsilon}(x^\star))} V_p^{\frac{1}{p}}(0) \right] \leq  2 \left[ \left( \frac{2q_\varepsilon}{\zeta}\right)^{\frac{1}{\mu}} +   \frac{e^{-\beta q_\varepsilon}}{\rho_t(\mathcal{B}_{r_\varepsilon}(x^\star))} V_p^{\frac{1}{p}}(0) \right]\leq \\
			&  2 \left[ \left( \frac{c \varepsilon^{\frac{1}{p}}  }{2}\right)^{\frac{1}{\mu}} +   \frac{e^{-\beta q_\varepsilon}}{\rho_t(\mathcal{B}_{r_\varepsilon}(x^\star))} V_p^{\frac{1}{p}}(0) \right]\leq \frac{c}{2}  V_p^{\frac{1}{p}}(0) + \frac{c m_\star}{2 m_\star} V_p^{\frac{1}{p}}(0)= K(0),
		\end{split}
	\end{equation*}
	where we used the definition of $C_{r_\varepsilon}$ in \eqref{eq:def_quantitative_laplace}, $q_\varepsilon$ and $\beta$ in \eqref{eq:beta}, and the mass bound in \eqref{eq:bound_mass}. Next, we show that $V_p(t)$ decays essentially exponentially fast in time. More precisely, we prove that $V_p(t)$ decays 
	\begin{enumerate}
		\item at least exponentially fast with rate $(1-\theta)C_{p,\alpha}$;
		\item at most exponentially fast with rate $(1+\theta/2)C_{p,\alpha}$,
	\end{enumerate}
	with $C_{p,\alpha}$ as in \eqref{eq:Cpalpha}. By definition of $T_\beta$,
	\begin{equation*}
		\begin{split}
			&V_{p-\alpha}(t) \vert \bar{x}_\mathcal{E}(t)-x^\star \vert^\alpha \leq G(d,p,\alpha) (V_p(t))^{\frac{p-\alpha}{p}} c^\alpha(V_p(t))^{\frac{\alpha}{p}} =  G(d,p,\alpha) c^\alpha V_p(t),  \\
			&(V_p(t))^{\frac{p-1}{p}} \vert \bar{x}_\mathcal{E}(t)-x^\star \vert \leq c (V_p(t))^{\frac{1}{p}} (V_p(t))^{\frac{p-1}{p}} = c V_p(t),
		\end{split}
	\end{equation*}
	with $c$ as in \eqref{eq:c} and $G(d,p,\alpha)$ as in \eqref{eq:G}.  By Proposition \ref{prop:proposition1} and by definition of $c$ in \eqref{eq:c}, we get  
	\begin{equation*}
		\frac{d}{dt} V_p(t) \leq -C_{p,\alpha} V_p(t) + \frac{\theta}{2} C_{p,\alpha}  V_p(t) + \frac{\theta}{2} C_{p,\alpha}  V_p(t) = -(1-\theta) C_{p,\alpha}  V_p(t). 
	\end{equation*}
	By Proposition \ref{prop:proposition2} we obtain a lower bound
	\begin{equation*}
		\frac{d}{dt} V_p(t)\geq -C_{p,\alpha} V_p(t) -\frac{\theta}{2}C_{p,\alpha} V_p(t) = -\left(1+\frac{\theta}{2}\right) C_{p,\alpha} V_p(t).
	\end{equation*}
	By Gr{\"o}nwall inequality we get 
	\begin{equation*}
		\begin{split}
			V_p(t) \leq  V_p(0) \exp\left(  -(1-\theta)C_{p,\alpha}t\right),\qquad
			V_p(t) \geq  V_p(0) \exp\left(  -\left(1+\frac{\theta}{2}\right)C_{p,\alpha}t\right).
		\end{split}
	\end{equation*}
	We now show that \eqref{eq:proof_goal} is satisfied. 
	\paragraph{Case $T_\beta \geq T_\star$.} 
	We have 
	\begin{equation*}
		V_p(T_\star) \leq V_p(0) \exp(-(1-\theta)C_{p,\alpha} T_\star) = V_p(0) \frac{\varepsilon}{V_p(0)} = \varepsilon.
	\end{equation*}
	Hence, by definition of $T_\star$, and by continuity of $V_p(t)$ we conclude that $V_p(T_\beta) =\varepsilon$ with $T_\beta = T_\star$. 
	\paragraph{Case $T_\beta <T_\star$ and $V_p(T_\star) \leq  \varepsilon$.} By continuity of $V_p(t)$, it holds that $V_p(T_\beta)=\varepsilon$. Thus,
	\begin{equation*}
		\varepsilon = V_p(T_\beta) \geq V_p(0) \exp\left(  -\left(1+\frac{\theta}{2}\right) C_{p,\alpha} T_\beta \right),
	\end{equation*}
	that is 
	\begin{equation*}
		\left( 1+\frac{\theta}{2} \right) T_\beta \geq \log\left( \frac{V_p(0)}{\varepsilon}\right)C_{p,\alpha}. 
	\end{equation*}
	Hence
	\begin{equation*}
		\frac{1-\theta}{1+\frac{\theta}{2}} T_\star= 	\frac{1}{(1+\frac{\theta}{2})C_{p,\alpha}} \log\left( \frac{V_p(0)}{\varepsilon}\right)C_{p,\alpha}\leq T_\beta < T_\star. 
	\end{equation*}
	\paragraph{Case $T_\beta <T_\star$ and $V_p(T_\star) > \varepsilon$.} We show that this case never occur due to the choice of $\beta$. More specifically, we verify that 
	\begin{equation}\label{eq:proof2} 
		\vert \bar{x}_\mathcal{E}(t)-x^\star \vert <K(T_\beta). 
	\end{equation}
	Indeed, fulfilling simultaneously both \eqref{eq:proof2} and $V_p(T_\beta )>\varepsilon$ would contradict the definition of $T_\beta$. By assumptions we get   
	\begin{equation*}
		\begin{split}
			\vert \bar{x}_\mathcal{E}(T_\beta)-x^\star \vert \leq & 2 \left[ C_{r_\varepsilon} +   \frac{e^{-\beta q_\varepsilon}}{\rho_{T_\beta}(\mathcal{B}_{r_\varepsilon}(x^\star))} V_p^{\frac{1}{p}}(T_\beta) \right]\leq\\ & \leq   \frac{c}{2}  V_p^{\frac{1}{p}}(T_\beta) + \frac{c m_\star}{2 m_\star} V_p^{\frac{1}{p}}(T_\beta)= c  V_p^{\frac{1}{p}}(T_\beta) = K(T_\beta). 
		\end{split}
	\end{equation*}
	This establishes the desired contradiction.
\end{proof}

\begin{remark}
	Note that Proposition \ref{prop:proposition1}-\ref{prop:proposition2} and Theorem \ref{prop:convergence} could be extended to the case with both classical and fractional diffusions, by assuming $\sigma>0$. 
\end{remark}

\section{Numerical methods }\label{sec:numerical_methods}
We now focus on the solution of the dynamics, relying on splitting approaches. Specifically, in each time interval $[0,\Delta t]$ we solve separately and consecutively the following problems 
\begin{equation}\label{eq:subpbs}
	\begin{cases}
		&\partial_t f^* = \mathcal{A}_{drift}(f^*),\\
		& f^*(x,0) = f(x,0) = f_0(x),
	\end{cases}\qquad 
	\begin{cases}
		&\partial_t f^{**}= \mathcal{A}_{diff}(f^{**}),\\
		& f^{**}(x,0) = f^*(x,\Delta t),
	\end{cases}\qquad 
\end{equation}
being $\mathcal{A}_{drift}$ and $\mathcal{A}_{diff}$ operators describing the drift and diffusion parts of equation \eqref{eq:fFP}. Let us denote by 
\begin{equation}
	\begin{split}
		&\mathcal{S}_1^{\Delta t} (f_0(x)) = f^*(x,\Delta t),\\ 
		&\mathcal{S}_2^{\Delta t} (f^{*}(x,\Delta t)) = f^{**}(x,\Delta t), 
	\end{split}
\end{equation}
the solution of the two subproblems in \eqref{eq:subpbs}, then we can write in compact notation
\begin{equation*}
	f^{**}(x,\Delta t) =    \mathcal{S}_2^{\Delta t}(\mathcal{S}_1^{\Delta t}(f_0(x)).
\end{equation*}
The method produces a first-order approximation of the true solution in the time interval $\Delta t$.

\subsection{A Nanbu-type scheme}
In order to approximate the evolution of the density function $f(x,t)$ we sample $N_s$ particles $x_i$, $i=1,\ldots,N_s$ from the initial distribution $f_0(x)$. We consider a time interval $[0,T]$ discretized in $N_t$ intervals of size $\Delta t$. We denote by $f^n$ the approximation of $f(x,n\Delta t)$ at time $t^n$. The next iterate is given by 
\begin{equation}\label{eq:f_splitting}
	\begin{cases}
		&f^* = f^n (1-\Delta t \lambda_1) + \Delta t \lambda_1 \mathcal{Q}_1^{+}(f^n),\\
		&f^{n+1} = f^* (1-\Delta t \lambda_2) + \Delta t \lambda_2 (\mathcal{Q}_2^{+}(f^*)+\mathcal{Q}_3^{+}(f^*)), 
	\end{cases}
\end{equation}
for $\lambda_1 \lambda_2 \geq 0$, and  
where, since $f^n$ is a probability density, mass conservation implies that
$\mathcal{Q}_1^{+}(f^n)$ is also a probability density. Consequently,
$f^*$ is a probability density, being a convex combination of probability densities. By the same argument, $f^{n+1}$ is a probability density. 
From a Monte Carlo viewpoint, we can interpreter \eqref{eq:f_splitting} as follows. With probability $(1-\Delta t \lambda_1)$ an agent in position $x$ will not interact with other agents, while with probability $\Delta t \lambda_1$ it will interact according to the interactions given by $\mathcal{Q}_1^+(f^n)$, and hence it will modify its position according to the drift, moving to position $x_*$. Subsequently, the agent $x_*$ with probability $(1-\Delta t \lambda_2)$ will not interact with other individuals, while with probability $\Delta t \lambda_2$ it will modify its position according the diffusion laws.  \\ To maximize the number of interactions, we assume $\lambda_1=\lambda_2 = 1/\Delta t$. Hence, for any $i=1\ldots,N_s$ we get
\begin{equation}\label{eq:binary_rules}
	\begin{cases}
		&x_i^* = x_i^n + \nu \Delta t (\bar{x}_\mathcal{E}^n - x_i^n),\\
		&x_i^{**} = x_i^* + \sigma \sqrt{\Delta t} D(\bar{x}_\mathcal{E},x_i^*) z_i+ \gamma (\Delta t)^{\frac{1}{\alpha}} D(\bar{x}_\mathcal{E},x_i^{*}) \tilde{z}_i, 
	\end{cases}
\end{equation}
with $\nu,\sigma,\gamma >0$, $z_i\sim \mathcal{N}(0,1)$ and $\tilde{z}_i\sim \mathcal{L}_t^\alpha$. 

The details of the numerical scheme are summarized in Algorithm \ref{alg:KBO}. The initial particle positions are sampled from the initial density, typically chosen as a uniform distribution in space unless prior information on the location of the global minimizer is available. The parameters $\delta_{stall}$ and $j_{stall}$ are introduced to detect convergence of the method. In particular, the algorithm is terminated if the distance between the consensus point $\bar{x}_\mathcal{E}$ and the global minimizer remains below a prescribed tolerance $\delta_{stall}$ for at least $j_{stall}$ consecutive iterations. Otherwise, the iterations proceed until the maximum number $N_t$ is reached.
\\
\begin{algorithm}{KBO with fractional diffusion}\label{alg:KBO}
	\begin{enumerate}
		\item[\texttt 1.] Draw $x_i^0$ with $i=1,\dots,N_s$ from the initial distribution $f^0(x)$ and set $n=0$, $j=0$.
		\item[\texttt 2.] Compute $\bar{x}_\mathcal{E}^0$ as in equation \eqref{eq:x_tot}.
		\item[\texttt 3.] \texttt{while} $n<N_t$ \texttt{and} $j<j_{stall}$
		\begin{enumerate}
			\item \texttt{for} $i=1$ \texttt{to}  $N$
			\begin{itemize}
				\item [] compute the new positions as in \eqref{eq:binary_rules}.
			\end{itemize}
			\texttt{end for}
			\item Compute $\bar{x}_\mathcal{E}^{n+1}$ as in equation \eqref{eq:x_tot}.
			\item  \texttt{if} $\Vert \bar{x}_\mathcal{E}^{n+1}-\bar{x}_\mathcal{E}^n\Vert_\infty\leq \delta_{stall}$
			\begin{enumerate}
				\item [] $j\leftarrow j+1$
			\end{enumerate}
			\texttt{end if} 
			\item [] $n\leftarrow n+1$
		\end{enumerate}
		\texttt{end while} 
	\end{enumerate}
\end{algorithm}
The above algorithm is inspired from Nanbu's method \cite{nanbu1980direct}. For a larger class of direct simulation Monte-Carlo algorithms for interacting particle dynamics, we refer to \cite{albi2013binary, pareschi2013interacting}. 

\section{Numerical experiments}\label{sec:numerical_experiments} 
We begin by presenting a validation test in a simplified scenario to assess the accuracy of the numerical method proposed in Section \ref{sec:numerical_methods}. We then proceed by testing the performance of the KBO method with a stable process $\alpha$ in terms of success rate and number of iterations, comparing its efficiency with the one of KBO with classical diffusion. 
\subsection{Validation test}
We will carry out this test in dimension one. Let us consider the simplified fractional Fokker-Planck equation
\begin{equation}\label{eq:FPE_validation}
	\partial_t f(x,t) = \partial_x (x f(x,t)) - (-\Delta)^{\alpha/2} (x^2 f(x,t)),
\end{equation}
coupled with the initial condition 
\begin{equation}\label{eq:FPE_validation_IC}
	f_0(x) = \frac{1}{\pi}\frac{1}{1+x^2}. 
\end{equation}
Let us notice that equation \eqref{eq:FPE_validation} shares many characteristics with equation \eqref{eq:fFP}.
In the case $\alpha = 1$, it can be proved that equation \eqref{eq:FPE_validation} with initial condition \eqref{eq:FPE_validation_IC} admits an exact solution given by 
\begin{equation}\label{eq:FPE_validation_exact}
	f_{ex}(x,t) = \frac{1}{\pi}\frac{\beta(t)}{\beta^2(t) + x^2}, \qquad \text{with} \quad \beta(t) = \frac{e^{-t}}{2-e^{-t}},
\end{equation}
for any $t\geq 0$. We consider a set of $N=10^6$ particles and we update the particle positions with the rules described in \eqref{eq:binary_rules}, assuming $\nu=1$, $\sigma=0$, $\bar{x}_\mathcal{E} = 0$, $D(\bar{x}_\mathcal{E},x) = x^2$, setting $\alpha = 1$. We let the dynamics to evolve up to time $T=2$, with time step $\Delta t = 0.01$. The density $f^N(x,t)$ is reconstructed over a one dimensional grid with $m_x=2^{10}$ points on the interval $[-20,20]$. In Figure \ref{fig:validation}, on the left, three snapshots of the dynamics at time $t=0.1$, $t=1$, and $t=2$.  The numerical and exact solution are in good agreement. On the right, the error between the exact and numerical solution at time $t=2$ computed as 
\begin{equation}\label{eq:validation_error}
	\mathcal{E}^N(T) = \Vert f_{ex}(x,T) - f^N(x,T) \Vert_\infty,
\end{equation}
for different values of $N$. As expected for a Monte Carlo method, the error decreases proportionally to $N^{-1/2}$. 
\begin{figure}[tbhp]
	\centering 
	\includegraphics[width=0.34\linewidth]{validation.eps}
	\includegraphics[width=0.35\linewidth]{validation_error_N.eps}
	\caption{ Validation test: on the left three snapshots of the numerical and exact solution as in \eqref{eq:FPE_validation_exact} for $N=10^6$ taken at time $t=0.1$, $t=1$, and $t=2$. On the right, the error as the number of particles $N$ increases between the exact solution in \eqref{eq:FPE_validation_exact} and its reconstruction, computed according to \eqref{eq:validation_error}.
	}
	\label{fig:validation}
\end{figure}


\subsection{KBO with fractional diffusion}
In this section, we investigate the efficiency of the KBO method driven by an $\alpha$-stable process. Unless otherwise specified, we consider the Rastrigin function
\begin{equation}\label{eq:Rastrigin}
	f(x) = 10 + \sum_{k=1}^d (x_k^2 - 10\cos(2\pi x_k)),
\end{equation}
as a benchmark problem, which admits a global minimizer at $x^\star = 0$.
We run $M=20$ simulations, and we consider one successful if 
\begin{equation}\label{eq:succ_estimate}
	\Vert \bar{x}_\mathcal{E}(t) - x^\star \Vert_\infty \leq 0.25,
\end{equation}
where $\Vert \cdot \Vert_\infty $ denotes the $L_\infty$ norm with respect to the dimension. 
We set $\beta = 5 \cdot 10^6$, and we adopt the trick described in \cite{fornasier2021consensus} to allow arbitrary large values of $\beta$. We set $N=200$ and assume agents are initially distributed in the hypercube $[-5.12,-2]^d$, which does not contain the global minimum. 
We let the dynamics to evolve for $N_t = 10^4$ iterations, with $\Delta t=0.1$. We set $j_{stall} = 10^3$, $\delta_{stall} = 10^{-4}$. In all tests, we set $\nu = 1$ and $\alpha = 1.5$. The other parameters, including the dimension $d$, will change in the different tests and will be specified later. 
\paragraph{Qualitative test.} We fix $d=2$, and we consider the non-convex non-differentiable function modified Alpine defined as:
\begin{equation}\label{eq:mod_alpine}
	f(x) = \sum_{k=1}^d(\abs{x_k \sin(x_k)}+0.2\abs{x_k}).
\end{equation}We consider both the case with purely classical diffusion, corresponding to $\sigma = 1$ and $\gamma = 0$, and the case combining classical and fractional diffusion, with $\sigma = \gamma = 1$. In dimension $d = 2$, the regime with only jumps ($\sigma = 0$, $\gamma = 1$) exhibits a behavior similar to the mixed diffusion case, and is therefore omitted for brevity.
In Figure \ref{fig:alpine_dynamics}, we report three snapshots at iterations $it = 10, 25, 150$ for both the classical KBO dynamics and the KBO dynamics with classical and fractional diffusion. In the latter case, particles exhibit a stronger exploratory behavior and progressively concentrate around the global minimizer. In contrast, in the purely classical setting, they tend to become trapped in a local minimum.
\begin{figure}[tbhp]
	\centering
	\includegraphics[width=0.327\linewidth]{mod_alpine_it10.eps}
	\includegraphics[width=0.327\linewidth]{mod_alpine_it25.eps}
	\includegraphics[width=0.327\linewidth]{mod_alpine_it150.eps}
	\caption{ Modified Alpine function: three snapshots of the KBO evolutions taken at iterations $it=10, 25, 150$ for fixed $d=2$, $\nu=1$. We set either $\sigma=1$ and $\gamma=0$ or $\sigma=1$ and $\gamma =1$ to model the case with pure classical diffusion and the one with both fractional and classical diffusion.  }
	\label{fig:alpine_dynamics}
\end{figure}
\paragraph{Test 1: Comparison in dimension $d=20$ for varying $\gamma$.}
We fix $d = 20$, and let $\gamma$ to vary between $\gamma = 1$ to $\sigma = 5$, testing the method on the Rastrigin function. We consider the KBO with pure $\alpha$-stable process setting $\sigma = 0$ and $\sigma = 3$. In Figure \ref{fig:rastr_gammarange} on the left, the plot of the success rate as $\gamma$ varies, and on the right, the plot of the iterations number. The method proved to be particularly effective, especially in the pure-jumps case ($\sigma = 0$) with $\gamma = 2.5$, where it achieves a higher success rate while requiring fewer iterations.
\begin{figure}[tbhp]
	\centering
	\includegraphics[width=0.327\linewidth]{rastr_succ_rate_gamma_splitting.eps}
	\includegraphics[width=0.327\linewidth]{rastr_iter_gamma_splitting.eps}
	\caption{ Success rate and mean number of iterations for the Rastrigin function for fixed $d=20$, $\nu=1$ as the parameter $\gamma$ varies. We set either $\sigma=0$ or $\sigma=3$ to model the case with a pure $\alpha$-stable process and the one with both fractional and classical diffusion.  Markers denote the value of the success rate and the number of iterations.
	}
	\label{fig:rastr_gammarange}
\end{figure}

\paragraph{Test 2: Comparison for varying dimension $d$.}
We let the dimension $d$ vary between $d=1$ and $d=50$, and we make different comparisons changing the values of $\sigma$ and $\gamma$, and testing the methods on the Rastrigin function. In Figure \ref{fig:rastr_drange} on the left, the plot of the success rate as $d$ varies, and on the right, the plot of the number of iterations. The KBO variant incorporating both an $\alpha$-stable process and diffusion ($\gamma=2$ and $\sigma=3$) demonstrates the highest effectiveness in terms of success rate. In contrast, the pure diffusion case ($\gamma=0$ and $\sigma=3$) achieves convergence in fewer iterations but exhibits a smaller success rate area. The best compromise is provided by the case with only the $\alpha$-stable process ($\gamma=2$ and $\sigma=0$), which combines a large success rate area with a reduced number of iterations.  For this choice of the parameters, the KBO with $\alpha$-stable process proves to be more effective than the classical KBO.   
\begin{figure}[tbhp]
	\centering
	\includegraphics[width=0.327\linewidth]{rastr_succ_rate_d20_splitting.eps}
	\includegraphics[width=0.327\linewidth]{rastr_iter_d20_splitting.eps} 
	\caption{Success rate and mean number of iterations for the Rastrigin function as the dimension $d$ varies. We consider pure $\alpha$-stable process setting $\gamma =2$ and $\sigma = 0$, pure diffusion setting $\gamma=0$ and $\sigma = 3$, and both diffusion and $\alpha$-stable process, setting $\gamma =2$ and $\sigma = 3$. Markers denote the value of the success rate and the number of iterations.
	}
	\label{fig:rastr_drange}
\end{figure}

\paragraph{Test 3: Comparison in dimension $d=20$ for varying $\sigma$.}
We fix $d = 20$, and let $\sigma$ to vary between $\sigma = 0$ to $\sigma = 6$, testing the method on the Rastrigin function. We consider the KBO with pure diffusion setting $\gamma = 0$, and the case with both diffusion and $\alpha$-stable process, setting $\gamma = 2$. In Figure \ref{fig:rastr_sigmarange} on the left, the plot of the success rate as $\sigma$ varies, and on the right, the plot of the number of iterations. The KBO with both diffusion and $\alpha$-stable process proved to be more effective than the one with only diffusion, showing a larger success rate area. 
\begin{figure}[tbhp]
	\centering
	\includegraphics[width=0.327\linewidth]{rastr_succ_rate_sigma_splitting.eps}
	\includegraphics[width=0.327\linewidth]{rastr_iter_sigma_splitting.eps} 
	\caption{Success rate and mean number of iterations for the Rastrigin function for fixed $d=20$, $\nu=1$ as the parameter $\sigma$ varies. The KBO with pure diffusion process is considered, setting $\gamma =0$, and the one with both diffusion and $\alpha$-stable, setting $\gamma = 2$.  Markers denote the value of the success rate and the number of iterations.
	}
	\label{fig:rastr_sigmarange}
\end{figure}

\paragraph{Test 4: Comparison of different benchmark functions.}\label{sec:test7}
We fix $d = 20$ and $\nu = 1$. We consider different differentiable and convex benchmark functions (see \cite{jamil2013literature} for a complete list), and we test the KBO method in the case with pure diffusion ($\gamma = 0$ and $\sigma = 3$), in the case with pure $\alpha$-stable process ($\gamma = 2$ and $\sigma = 0$), and in the case with both $\alpha$-stable process and diffusion ($\gamma = 2$ and $\sigma = 3$).  
Figure \ref{fig:funct_in0} reports a comparison in terms of success rate and mean number of iterations. For the different parameter choices, the KBO method with pure classical diffusion, pure jumps, and combined jumps and diffusion achieves a higher success rate. The fractional KBO outperforms the classical version for the Rastrigin function, whereas for the Rosenbrock function, the KBO with classical diffusion attains a higher success rate. Moreover, methods employing either pure jumps or pure diffusion converge faster than the approach combining classical diffusion and jumps. 
\begin{figure}[tbhp]
	\centering
	\includegraphics[width=0.327\linewidth]{funct_succ_rate_splitting.eps}
	\includegraphics[width=0.327\linewidth]{funct_iter_splitting.eps} 
	\caption{ Success rate and mean number of iterations for the different benchmark functions for fixed $d=20$.  We consider pure diffusion ($\gamma = 0$, and $\sigma = 3$), pure  $\alpha$-stable process ($\gamma = 2$ and $\sigma = 0$), and both $\alpha$-stable process and diffusion ($\gamma = 2$ and $\sigma = 3$). Markers denote the value of the success rate and the number of iterations for the different benchmark functions.
	}
	\label{fig:funct_in0}
\end{figure}
We now test the method on several non-differentiable functions, some of which are also non-convex (see \cite{karmitsa2007test} for a complete list). Since the method is gradient-free, convergence is expected. This is confirmed by the results shown in Figure \ref{fig:funct_nondiff}, where the success rate (left) and the number of iterations to convergence (right) are reported. All three variants exhibit a high success-rate region and a comparable number of iterations required to reach convergence.
\begin{figure}[tbhp]
	\centering
	\includegraphics[width=0.327\linewidth]{funct_succ_rate_splitting_non_diff.eps}
	\includegraphics[width=0.327\linewidth]{funct_iter_splitting_non_diff.eps} 
	\caption{Success rate and mean number of iterations for the different benchmark functions for fixed $d=20$.  We consider pure diffusion ($\gamma = 0$, and $\sigma = 3$), pure $\alpha$-stable process ($\gamma = 2$ and $\sigma = 0$), and both $\alpha$-stable process and diffusion ($\gamma = 2$ and $\sigma = 3$). Markers denote the value of the success rate and the number of iterations for the different benchmark functions.
	}
	\label{fig:funct_nondiff}
\end{figure}


\section{Conclusions}
\label{sec:conclusions}

In this work, we have introduced a novel variant of the classical KBO method based on fractional diffusion driven by an $\alpha$-stable stochastic process. By incorporating nonlocal stochastic effects, the proposed approach enhances the exploration of complex, high-dimensional, and non-convex landscapes. Starting from the particle-level formulation, we derived the corresponding fractional Fokker–Planck equation via a Fourier representation and provided a rigorous convergence analysis, thereby offering both modeling and theoretical support for the method.
The numerical experiments confirm that the fractional diffusion mechanism can improve performance with respect to standard diffusion-based approaches, particularly in challenging multimodal scenarios, where avoiding premature convergence is crucial.
Several directions for future research remain open. In particular, a promising perspective is the application of the proposed framework to the training of neural networks, where the optimization landscape is notoriously high-dimensional and non-convex. The nonlocal exploratory features of the method could provide advantages in escaping poor local minima and improving generalization. Further investigations will also focus on extending the method to multi-objective optimization problems.


\section*{Acknowledgments}
This work has been written within the activities of GNCS group of INdAM (Italian National Institute of High Mathematics). The research of FF was funded with the contribution of the Italian Ministry of University and Research (MUR) – Fondo Italiano per la Scienza Bando FIS2 - CUP F53C25000210001 – project code FIS-2023-0133. The research of MH has been supported by the Deutsche Forschungsgemeinschaft (German Research Foundation) through HE5386/33-1 Control of Interacting Particle Systems, and Their Mean-Field, and Fluid-Dynamic Limits (560288187) and 
HE5386/34-1 Partikelmethoden für unendlich dimensionale Optimierung (561130572). PAS acknowledges support from the Simons Foundation through the Travel Support for Mathematicians program.  GA acknowledges support by MUR project PRIN 2022 PNRR nr.~P2022JC95T, CUP: B53D23027840001.

\bibliographystyle{abbrv}
\bibliography{references}

\end{document}
